\ifdefined\commandsFileStandaloneCompilation
\else
\documentclass [11pt,twoside]{article}
\usepackage[utf8]{inputenc}
\usepackage[T1]{fontenc}


%Page margins, header and footer positions
\usepackage{geometry}
 \geometry{
 a4paper,
 total={210mm,297mm},
 left=25mm,
 right=25mm,
 top=30mm,
 bottom=25mm,
 headsep=7mm}

\interfootnotelinepenalty=10000

%To display filling dots in the TOC for all entries
\usepackage[titles]{tocloft}
\renewcommand{\cftsecleader}{\cftdotfill{\cftdotsep}}

%Define new header and footer style
\usepackage{fancyhdr}

\pagestyle{fancy}
\fancyhf{}
\lhead{\color{Gray}{\small{Best Bike Paths Project}}}
\lfoot{\textcolor{Gray}{\small{Copyright © 2025, A. Ferritti, J. Lin, T. Tosca – All rights reserved}}}
\rfoot{\textcolor{Gray}{\thepage}}
\renewcommand{\headrulewidth}{0pt}

%PACKAGES
\usepackage{wasysym}
\usepackage{pifont}

\newcommand{\supported}{\ding{52}\xspace}
\newcommand{\unsupported}{\ding{55}\xspace}
\newcommand{\partsupported}{\textcolor{black!40}{\ding{52}}\xspace}
\newcommand{\lowsupported}{\textcolor{black!20}{\ding{52}}\xspace}
\newcommand{\unknowsupported}{\textbf{?}\xspace}

%Font: Times
\usepackage{times}
%Change monospaced font
\renewcommand{\ttdefault}{lmtt}

%tables
\usepackage{tabu}
\usepackage{tabularx}
\usepackage{ltablex}
\usepackage{longtable}
\usepackage{float} % To allow the use of H modifier in long tables

%landscape mode
\usepackage{pdflscape}
\usepackage{rotating}
\usepackage{caption}

%make landscape mode be sensitive to even and odd pages
%start
\def\myrotate{\ifodd\c@page\else-\fi 90}
\makeatletter
\global\let\orig@begin@landscape=\landscape%
\global\let\orig@end@landscape=\endlandscape%
\gdef\@true{1}
\gdef\@false{0}
\gdef\landscape{%
    \global\let\within@landscape=\@true%
    \orig@begin@landscape%
}%
\gdef\endlandscape{%
    \orig@end@landscape%
    \global\let\within@landscape=\@false%
}%
\@ifpackageloaded{pdflscape}{%
    \gdef\pdf@landscape@rotate{\PLS@Rotate}%
}{
    \gdef\pdf@landscape@rotate#1{}%
}
\let\latex@outputpage\@outputpage
\def\@outputpage{
    \ifx\within@landscape\@true%
        \if@twoside%
            \ifodd\c@page%
                \gdef\LS@rot{\setbox\@outputbox\vbox{%
                    \pdf@landscape@rotate{-90}%
                    \hbox{\rotatebox{90}{\hbox{\rotatebox{180}{\box\@outputbox}}}}}%
                }%
            \else%
                \gdef\LS@rot{\setbox\@outputbox\vbox{%
                    \pdf@landscape@rotate{+90}%
                    \hbox{\rotatebox{90}{\hbox{\rotatebox{0}{\box\@outputbox}}}}}%
                }%
            \fi%
        \else%
            \gdef\LS@rot{\setbox\@outputbox\vbox{%
                \pdf@landscape@rotate{+90}%
                \hbox{\rotatebox{90}{\hbox{\rotatebox{0}{\box\@outputbox}}}}}%
            }%
        \fi%
    \fi%
    \latex@outputpage%
}
\makeatother
%end

%graphics
\usepackage{graphicx}
\usepackage[dvipsnames, table]{xcolor}
%If you upload images from PC, you need to insert code for the path here (different for Windows and Unix OS)

%References
%\usepackage{xpatch}
%\usepackage[backend=biber, style=numeric, citestyle=numeric, sorting=none]{biblatex}
%\addbibresource{main.bib}

%Other
\usepackage{ifthen}
\usepackage{xspace}
\usepackage{enumitem}
\usepackage{amssymb}
\usepackage[pdftex, colorlinks, linkcolor=Blue, citecolor=Blue, urlcolor=Blue]{hyperref}
\newcommand{\comment}[1]{{\color{Red}$\blacktriangleright$ Comment: #1 $\blacktriangleleft$}}


% Some utilities\ldots
\usepackage{soul}
\usepackage{tikz}

\usetikzlibrary{calc}
\usetikzlibrary{decorations.pathmorphing}


\makeatletter

\newcommand{\defhighlighter}[3][]{%
  \tikzset{every highlighter/.style={color=#2, fill opacity=#3, #1}}%
}

\defhighlighter{yellow}{.5}

\newcommand{\highlight@DoHighlight}{
  \fill [ decoration = {random steps, amplitude=1pt, segment length=15pt}
        , outer sep = -15pt, inner sep = 0pt, decorate
       , every highlighter, this highlighter ]
        ($(begin highlight)+(0,8pt)$) rectangle ($(end highlight)+(0,-3pt)$) ;
}

\newcommand{\highlight@BeginHighlight}{
  \coordinate (begin highlight) at (0,0) ;
}

\newcommand{\highlight@EndHighlight}{
  \coordinate (end highlight) at (0,0) ;
}

\newdimen\highlight@previous
\newdimen\highlight@current

\DeclareRobustCommand*\highlight[1][]{%
  \tikzset{this highlighter/.style={#1}}%
  \SOUL@setup
  %
  \def\SOUL@preamble{%
    \begin{tikzpicture}[overlay, remember picture]
      \highlight@BeginHighlight
      \highlight@EndHighlight
    \end{tikzpicture}%
  }%
  %
  \def\SOUL@postamble{%
    \begin{tikzpicture}[overlay, remember picture]
      \highlight@EndHighlight
      \highlight@DoHighlight
    \end{tikzpicture}%
  }%
  %
  \def\SOUL@everyhyphen{%
    \discretionary{%
      \SOUL@setkern\SOUL@hyphkern
      \SOUL@sethyphenchar
      \tikz[overlay, remember picture] \highlight@EndHighlight ;%
    }{%
    }{%
      \SOUL@setkern\SOUL@charkern
    }%
  }%
  %
  \def\SOUL@everyexhyphen##1{%
    \SOUL@setkern\SOUL@hyphkern
    \hbox{##1}%
    \discretionary{%
      \tikz[overlay, remember picture] \highlight@EndHighlight ;%
    }{%
    }{%
      \SOUL@setkern\SOUL@charkern
    }%
  }%
  %
  \def\SOUL@everysyllable{%
    \begin{tikzpicture}[overlay, remember picture]
      \path let \p0 = (begin highlight), \p1 = (0,0) in \pgfextra
        \global\highlight@previous=\y0
        \global\highlight@current =\y1
      \endpgfextra (0,0) ;
      \ifdim\highlight@current < \highlight@previous
        \highlight@DoHighlight
        \highlight@BeginHighlight
      \fi
    \end{tikzpicture}%
    \the\SOUL@syllable
    \tikz[overlay, remember picture] \highlight@EndHighlight ;%
  }%
  \SOUL@
}

\makeatother

% Common abbrev. are set as commands to ensure proper spacing after the dot
\RequirePackage{xspace}
\newcommand{\ie}{i.e.\@\xspace}
\newcommand{\aka}{a.k.a.\@\xspace}
\newcommand{\Ie}{I.e.\@\xspace}
\newcommand{\cf}{cf.\@\xspace}
\newcommand{\Cf}{Cf.\@\xspace}
\newcommand{\eg}{e.g.\@\xspace}
\newcommand{\Eg}{E.g.\@\xspace}
\newcommand{\etal}{et al.\@\xspace}
\newcommand{\etc}{etc.\@\xspace}
\newcommand{\wrt}{w.r.t.\@\xspace}
\newcommand{\Wrt}{W.r.t.\@\xspace}



\date{}

% Detect whether this file is being compiled directly rather than being
% \input{} from main.tex so we can bootstrap the document correctly.
\newif\ifcommandsFileStandalone
\commandsFileStandalonefalse
\ifcsname pdfstrcmp\endcsname
    \ifnum\pdfstrcmp{\jobname}{commandsFile}=0\relax
        \commandsFileStandalonetrue
    \fi
\fi

\ifcommandsFileStandalone
    \newcommand{\commandsFileStandaloneCompilation}{}
    \def\commandsFileStandaloneBootstrap{\ifdefined\commandsFileStandaloneCompilation
\else
\documentclass [11pt,twoside]{article}
\usepackage[utf8]{inputenc}
\usepackage[T1]{fontenc}


%Page margins, header and footer positions
\usepackage{geometry}
 \geometry{
 a4paper,
 total={210mm,297mm},
 left=25mm,
 right=25mm,
 top=30mm,
 bottom=25mm,
 headsep=7mm}

\interfootnotelinepenalty=10000

%To display filling dots in the TOC for all entries
\usepackage[titles]{tocloft}
\renewcommand{\cftsecleader}{\cftdotfill{\cftdotsep}}

%Define new header and footer style
\usepackage{fancyhdr}

\pagestyle{fancy}
\fancyhf{}
\lhead{\color{Gray}{\small{Best Bike Paths Project}}}
\lfoot{\textcolor{Gray}{\small{Copyright © 2025, A. Ferritti, J. Lin, T. Tosca – All rights reserved}}}
\rfoot{\textcolor{Gray}{\thepage}}
\renewcommand{\headrulewidth}{0pt}

%PACKAGES
\usepackage{wasysym}
\usepackage{pifont}

\newcommand{\supported}{\ding{52}\xspace}
\newcommand{\unsupported}{\ding{55}\xspace}
\newcommand{\partsupported}{\textcolor{black!40}{\ding{52}}\xspace}
\newcommand{\lowsupported}{\textcolor{black!20}{\ding{52}}\xspace}
\newcommand{\unknowsupported}{\textbf{?}\xspace}

%Font: Times
\usepackage{times}
%Change monospaced font
\renewcommand{\ttdefault}{lmtt}

%tables
\usepackage{tabu}
\usepackage{tabularx}
\usepackage{ltablex}
\usepackage{longtable}
\usepackage{float} % To allow the use of H modifier in long tables

%landscape mode
\usepackage{pdflscape}
\usepackage{rotating}
\usepackage{caption}

%make landscape mode be sensitive to even and odd pages
%start
\def\myrotate{\ifodd\c@page\else-\fi 90}
\makeatletter
\global\let\orig@begin@landscape=\landscape%
\global\let\orig@end@landscape=\endlandscape%
\gdef\@true{1}
\gdef\@false{0}
\gdef\landscape{%
    \global\let\within@landscape=\@true%
    \orig@begin@landscape%
}%
\gdef\endlandscape{%
    \orig@end@landscape%
    \global\let\within@landscape=\@false%
}%
\@ifpackageloaded{pdflscape}{%
    \gdef\pdf@landscape@rotate{\PLS@Rotate}%
}{
    \gdef\pdf@landscape@rotate#1{}%
}
\let\latex@outputpage\@outputpage
\def\@outputpage{
    \ifx\within@landscape\@true%
        \if@twoside%
            \ifodd\c@page%
                \gdef\LS@rot{\setbox\@outputbox\vbox{%
                    \pdf@landscape@rotate{-90}%
                    \hbox{\rotatebox{90}{\hbox{\rotatebox{180}{\box\@outputbox}}}}}%
                }%
            \else%
                \gdef\LS@rot{\setbox\@outputbox\vbox{%
                    \pdf@landscape@rotate{+90}%
                    \hbox{\rotatebox{90}{\hbox{\rotatebox{0}{\box\@outputbox}}}}}%
                }%
            \fi%
        \else%
            \gdef\LS@rot{\setbox\@outputbox\vbox{%
                \pdf@landscape@rotate{+90}%
                \hbox{\rotatebox{90}{\hbox{\rotatebox{0}{\box\@outputbox}}}}}%
            }%
        \fi%
    \fi%
    \latex@outputpage%
}
\makeatother
%end

%graphics
\usepackage{graphicx}
\usepackage[dvipsnames, table]{xcolor}
%If you upload images from PC, you need to insert code for the path here (different for Windows and Unix OS)

%References
%\usepackage{xpatch}
%\usepackage[backend=biber, style=numeric, citestyle=numeric, sorting=none]{biblatex}
%\addbibresource{main.bib}

%Other
\usepackage{ifthen}
\usepackage{xspace}
\usepackage{enumitem}
\usepackage{amssymb}
\usepackage[pdftex, colorlinks, linkcolor=Blue, citecolor=Blue, urlcolor=Blue]{hyperref}
\newcommand{\comment}[1]{{\color{Red}$\blacktriangleright$ Comment: #1 $\blacktriangleleft$}}


% Some utilities\ldots
\input{util/highlight}
\input{util/abbrev}



\date{}

% Detect whether this file is being compiled directly rather than being
% \input{} from main.tex so we can bootstrap the document correctly.
\newif\ifcommandsFileStandalone
\commandsFileStandalonefalse
\ifcsname pdfstrcmp\endcsname
    \ifnum\pdfstrcmp{\jobname}{commandsFile}=0\relax
        \commandsFileStandalonetrue
    \fi
\fi

\ifcommandsFileStandalone
    \newcommand{\commandsFileStandaloneCompilation}{}
    \def\commandsFileStandaloneBootstrap{\input{main}}
\else
    \let\commandsFileStandaloneBootstrap\relax
\fi
\commandsFileStandaloneBootstrap

\fi

\usepackage{tikz}

\usepackage{tabularx}
\usepackage[dvipsnames]{xcolor}
\usepackage{listings}
\usepackage{Alloy/alloy-style}


\usepackage[most]{tcolorbox}

\tcbset{
  chatstyle/.style={
    breakable,
    enhanced,
    rounded corners,    % angoli arrotondati
    arc=2mm,            % raggio arrotondamento
    boxrule=1pt,
    coltitle=black,      % Colore del testo del titolo
    fonttitle=\bfseries, % Titolo in grassetto
    attach title to upper,
    bottomtitle=1mm,
    toptitle=1mm,
    titlerule=0.5pt
  }
}

% Box per il Prompt
\newtcolorbox{PromptBox}{
  chatstyle,
  colback=blue!5,
  colframe=blue!60!black,
  title=Prompt
}

% Box per ChatGPT corretto (verde)
\newtcolorbox{ChatGPTGood}{
  chatstyle,
  colback=green!5,
  colframe=green!60!black,
  title=Gemini
}

% Box per ChatGPT errato (rosso)
\newtcolorbox{ChatGPTBad}{
  chatstyle,
  colback=red!5,
  colframe=red!60!black,
  title=Gemini
}

\begin{document}

%TITLE PAGE

\begin{titlepage}

% --- BLOCCO SEGNALIBRO ---
\begin{tikzpicture}[remember picture, overlay]
    \coordinate (TopLeft) at (current page.north west);
    
    % --- MODIFICA QUI ---
    % Sostituisci '\draw' con '\fill'
    \fill[Blue, opacity=0.8] % <<< CAMBIATO QUI
        ([xshift=3cm]TopLeft) 
        -- ++(0, -6cm)        
        -- ++(1.25cm, 1.5cm)  
        -- ++(1.25cm, -1.5cm) 
        -- ++(0, 6cm)         
        -- cycle;             
\end{tikzpicture}
% --- FINE BLOCCO SEGNALIBRO ---


% --- INTESTAZIONE: INFO PROGETTO E LOGO ---
\noindent 
\begin{minipage}[t]{0.7\textwidth} 
    % Vuoto
\end{minipage}
\hfill 
\begin{minipage}[t]{0.45\textwidth} 
    \vspace{0pt} 
    \flushright 
    \includegraphics[width=\linewidth]{Images/PolimiLogo}
\end{minipage}

\vfill 

% --- TITOLO E SOTTOTITOLO ---
\begin{flushleft}
    {\textcolor{Blue}{\textbf{\Huge{Requirement Analysis and \\ Specification Document}}}}
    
    \vspace{1cm} 
    {\textcolor{MyGray}{\textbf{\LARGE{Best Bike Paths Project}}}}
\end{flushleft} 

\vfill

% --- PIÈ DI PAGINA: DATA ---
\begin{flushleft} 
    {\large 2025-2026}
\end{flushleft} 

\end{titlepage}
%Define deliverable specific info
%Replace cell contents where needed
\begin{table}[h!]
\begin{tabu} to \textwidth { X[0.3,r,p] X[0.7,l,p] }
\hline

\textbf{Deliverable:} & RASD\\
\textbf{Title:} & Requirement Analysis and Specification Document \\
\textbf{Authors:} & Andrea Ferritti, Jieke Lin, Tommaso Tosca \\
\textbf{Version:} & 1.0 \\ 
\textbf{Date:} & 15-November-2025 \\
\textbf{Download page:} & https://github.com/ferritti/FerrittiLinTosca \\
\textbf{Copyright:} & Copyright © 2025, A. Ferritti, J. Lin, T. Tosca – All rights reserved \\
\hline
\end{tabu}
\end{table}




\setcounter{page}{2}


%------------------------------------------------------------------------------------------------------------------------------------------------
\newpage
\addcontentsline{toc}{section}{Table of Contents}
\tableofcontents
\newpage
\addcontentsline{toc}{section}{List of Figures}
\listoffigures
\addcontentsline{toc}{section}{List of Tables}
\listoftables

%------------------------------------------------------------------------------------------------------------------------------------------------
\clearpage
{\color{Blue}{\section{Introduction}}}
\label{sect:introduction}
\subsection{Purpose}
The growing popularity of cycling as both a means of transportation and a recreational activity has increased the need for reliable, up-to-date information about bike paths. Cyclists often struggle to identify safe and efficient routes, especially when facing incomplete or outdated data regarding path conditions, obstacles, or environmental factors. At the same time, users who wish to contribute to the cycling community by sharing information about the paths they travel may find it difficult to organize, store, and communicate such data in a structured and meaningful way.
The purpose of BBP is to provide an environment where cyclists can record their trips, contribute information about bike paths, and access curated, community-driven data to make informed decisions about their routes. Through a combination of manual input and automated data collection from mobile devices, users can keep track of their cycling activities while simultaneously improving the shared knowledge base about path conditions.
BBP aims to give all users—whether registered or not—a clear understanding of the quality and safety of available bike paths, offering route visualizations, condition assessments, and recommendations based on consolidated and verified information. Transparency, reliability, and community collaboration are promoted at every step, ensuring that cyclists can navigate efficiently while benefiting from accurate and continuously updated path information.
\\

\noindent
BBP, therefore, intends to serve the following primary goals:

\begin{description}
    \item[\textbf{[G1]}] Registered users should have access to performance metrics about their trips.
    
    \item[\textbf{[G2]}] When available, users would like to be informed about meteorological information related to the paths of interest.
    
    \item[\textbf{[G3]}] Registered users should have the possibility to insert informations about the path of interest.
    
    \item[\textbf{[G4]}] Information about bike paths provided by registered users should be acquirable both manually and automatically.
    
    \item[\textbf{[G5]}] Automatically detected issues should be validated by the registered user before becoming available to others.
    
    \item[\textbf{[G6]}] Registered users should be able to make their provided information publishable to the community.
    
    \item[\textbf{[G7]}] Users should be able to visualize viable bike paths between two points.
    
    \item[\textbf{[G8]}] Users should be able to compare alternative bike paths based on their overall score.
\end{description}


\subsection{Scope}
The Best Bike Paths (BBP) system aims to support the biking community by mediating the interactions between individual cyclists and the shared repository of information regarding bike paths. The external users who interact with the platform include registered cyclists, unregistered visitors, and potentially local cycling associations or municipalities interested in monitoring the conditions of bike routes. Since only registered users can contribute new information or validate automatically collected data, the system must identify these users through personal accounts, while unregistered visitors access the platform in a read-only mode. Registered cyclists provide BBP with details about their trips either manually or through an automated acquisition mode that collects GPS data and movement information. In addition to recording the path itself, the application also provides users with performance metrics derived from their trips, such as total distance, duration, average speed, and other relevant indicators to support personal activity monitoring and analysis. These data sources become key assets for enhancing trip analysis, improving the accuracy of anomaly detection, and increasing the overall reliability of the platform’s scoring mechanisms.
When available, the system also enriches trip data with meteorological information, such as weather conditions, temperature, and wind speed, retrieved from external services.
The system supports users in reviewing, correcting, and publishing the collected information, and it aggregates contributions from multiple cyclists to ensure that the status of each path reflects fresh and trustworthy evidence. All external users can request the visualization of one or more paths between an origin and a destination, and BBP assists them by computing ranked paths based on their status and effectiveness. Through these features, the application facilitates the discovery of safe and efficient bike paths and provides a continuously updated overview of the cycling network.



\subsubsection{World Phenomena}
\begin{description}
    \item[\textbf{[WP1]}] A user moves along a path.
    
    \item[\textbf{[WP2]}] The user’s physical speed varies over time.
    
    \item[\textbf{[WP3]}] The user’s device experiences vibrations or impacts due to road conditions.
    
    \item[\textbf{[WP4]}] The user physically reaches the starting location of the trip.
    
    \item[\textbf{[WP5]}] The user arrives at the ending location of the trip.
    
    \item[\textbf{[WP6]}] The user trip is conditioned by the weather conditions.
\end{description}

\subsubsection{Shared Phenomena}
\paragraph{World Controlled}
\begin{description}
    \item[\textbf{[SP1]}]  The user publishes information about the path in manual mode

    \item[\textbf{[SP2]}] The user accepts path information suggested by automated mode

    \item[\textbf{[SP3]}] The system receives accelerometer readings generated by vibrations

    \item[\textbf{[SP4]}] The system receives gyroscope readings generated by device rotation 

    \item[\textbf{[SP5]}] The system receives GPS location data from the cyclist’s device

    \item[\textbf{[SP6]}]  Different users may observe the same path and report different aspects of it

    \item[\textbf{[SP7]}] The user inserts an origin and a destination

    \item[\textbf{[SP8]}] The user starts the navigation
\end{description}


\paragraph{Machine Controlled}
\begin{description}
    \item[\textbf{[SP9]}] The system suggests path information to be inserted in automated mode 

    \item[\textbf{[SP10]}] The system provides the top 5 paths based on the score

    \item[\textbf{[SP11]}] The system provides information about the path of interest

    \item[\textbf{[SP12]}] The system provides performance metrics about the user’s trip
\end{description}
\subsection{Definitions, Acronyms, Abbreviations}
\subsubsection{Definitions}

\begin{itemize}
    \item \textbf{Bike Path}: A road segment or sequence of segments that is suitable for bicycles. A path is considered bike-safe if it includes a dedicated bike lane or if car traffic is rare and speed limits are compatible with typical bike speeds.
    \item \textbf{Trip}: A cycling activity recorded by a registered user. A trip includes GPS tracking data, performance metrics (distance, average speed, etc.), and, when available, meteorological information retrieved from an external weather service.
    \item \textbf{Performance Metrics}: Statistics associated with a trip, such as total distance, duration, average speed, maximum speed, elevation gain, and other cycling-related measures.
    \item \textbf{Meteorological Information}: External data enriched by BBP when available, including temperature, weather conditions (sunny, cloudy, rainy…), wind speed, humidity, and similar environmental indicators.
    \item \textbf{Manual Mode}: A mode in which a user manually enters information about a bike path, specifying street names, path segments, their status, and any obstacles.
    \item \textbf{Automated Mode}: A mode in which BBP automatically collects data from the user’s mobile device while biking. Data includes GPS traces to infer the following path, accelerometer and gyroscope readings to detect potential potholes or obstacles, and speed information to infer that the user is cycling.
    \item \textbf{Obstacle}: Any relevant physical irregularity affecting the bike path (e.g., potholes, bumps, debris). In automated mode, obstacles are inferred from device sensors and later confirmed by the user.
    \item \textbf{Status}: A qualitative assessment of a bike path segment. Possible values include: optimal, medium, sufficient, requires maintenance, etc. Status can be user-generated (manual or automatic) and may be aggregated from multiple users.
    \item \textbf{Information Publishable}: Verified and user-approved data about a bike path (status, obstacles) that can be made publicly accessible to the community.
    \item \textbf{Score}: A computed score used to rank alternative bike paths between an origin and a destination. The score considers both the quality (status) of the path and its effectiveness in connecting the two points.
\end{itemize}

\subsubsection{Acronyms}
\begin{itemize}
    \item \textbf{BBP:} Best Bike Path.
    \item \textbf{GPS:} Global Positioning System.
\end{itemize}

\subsubsection{Abbreviations}
\begin{itemize}
    \item \textbf{G:} Goal.
    \item \textbf{SP:} Shared Phenomena.
    \item \textbf{WP: } World Phenomena.
\end{itemize}

\subsection{Revision history}
Version 1.0: 23/12/2025
\subsection{Reference Documents}
\begin{itemize}
    \item Assignment RDD AY 2025-2026
    \item Software Engineering 2 AY 2025/2026 slides "CreatingRASD"
\end{itemize}
\subsection{Document Structure}
\begin{enumerate}
    \item \textbf{Introduction:} It aims to provide a brief overview of the project, focusing in particular on the motivations behind it and the goals intended to be achieved through its development.
    \item \textbf{Overall Description:} It is a high-level description of how the system operates, including a detailed explanation of the phenomena involving the world, the machine, or both. This section also provides the domain description along with the related assumptions.
    \item \textbf{Specific Requirements:} This section presents a detailed analysis of the requirements needed to achieve the identified goals. In addition, it provides further information relevant to developers, such as hardware and software constraints.
    \item \textbf{Formal analysis:} It provides a formal specification of the world phenomena using Alloy.
    \item \textbf{Effort spent:} It reports the time spent to produce this document, organized by section.
    \item \textbf{References:}It includes references to all documents and software used in the preparation of this paper.
\end{enumerate}
 

%------------------------------------------------------------------------------------------------------------------------------------------------
\clearpage
{\color{Blue}{\section{Overall Description}}}
\label{sect:overview}
\subsection{Product perspective}

\subsubsection{Scenarios}

    \textbf{Registered users} 
    \vspace{2mm}
    \noindent


    \noindent
    \textbf{User publishes path information in automated mode} 
    \vspace{1mm}
    
    \noindent
    After completing the bike ride,  user U clicks the stop button related to the automated mode and will be redirected to the path data entry page with all the field 
    pre-filled by the system based on the data it collected, with the possibility to modify the status and the obstacles. Finally, U presses the "Publish" button, and the new path becomes immediately visible to all users.

    \vspace{3mm}
    \noindent
    \textbf{User publishes path information in manual mode} 
    \vspace{1mm}
    
    \noindent
    User U wants to share a path they frequently use with the community. From the main map interface, U presses the "Create Path" button (the plus icon). U enters the drawing mode and taps a sequence of points on the map to define the itinerary; the application automatically interprets these inputs and snaps the route line to the correct street geometry. If U is aware of specific obstacles along this route, such as deep potholes or dangerous intersections, U switches the tool to "Obstacles" mode and taps on the corresponding points of the path to flag them. Once the geometry is defined, U proceeds to the next step. The application presents a summary form where U enters the estimated duration and selects a status rating. The system automatically lists the names of the streets detected along the drawn path. Finally, U presses the "Publish" button, and the new path becomes immediately visible to all users.

    \vspace{3mm}
    \noindent
    \textbf{User eliminates a published path} 
    \vspace{1mm}
    
    \noindent
    User U wants to delete a previously published route. In the navbar, U clicks on “Published,” where he will find all the paths they have published. U clicks on the route they want to delete; the page for editing the route’s information will open, and they click on the icon associated with deleting the route.

    
    \vspace{3mm}
    \noindent
    \textbf{User edits the published path}
    \vspace{1mm}
    
    \noindent
    User U views their list of shared routes in the "Published" section and decides to update the details of a specific path due to changed road conditions. U taps on the path card to open the detailed view. To make changes, U presses the "Edit" button (the pencil icon) located in the header. The interface transitions to editing mode: the Start and End location names become editable text fields, and the read-only status badge is replaced by a selection grid, allowing U to quickly switch the rating (e.g., from "Maintenance Required" back to "Optimal"). Additionally, the description field becomes an active text area, allowing U to refine the route notes. Once the modifications are complete, U presses the "Save" button. The application updates the record and immediately reverts to the read-only detailed view, displaying the refreshed information.

    \vspace{3mm}
    \noindent
    \textbf{User modifies personal information}
    \vspace{1mm}
    
    \noindent
    User U wants to edit their personal information. In the navbar, U clicks on “Profile” and then on “Settings”, where he will find all their personal information with the option to modify it

    \newpage
    \noindent
    \textbf{User logs into the application}
    \vspace{1mm}
    
    \noindent
    User U, currently viewing the Home Page in the unregistered state, clicks on "My Profile" in the navbar. This action directs them to the Login screen. On this screen, U is prompted to enter their credentials: they fill in the email address in the designated field and the corresponding password in the password field. After ensuring the data is correct, U clicks on the "Login" button. Upon successful validation, the system grants U access and redirects them back to the main application Home Page, now displaying features and content available to a registered user.
    

    \vspace{3mm}
    \noindent
    \textbf{User discards the collected path information}
    \vspace{1mm}
    
    \noindent
    User U wants to discard the path information he entered. On the data entry screen for the given route, U clicks the “Discard” button, and the information will not be published to the community. However, if the information was entered in automated mode, it will still be added to the recent trips.

    \vspace{3mm}
    \noindent
    \textbf{User saves the trip}
    \vspace{1mm}
    
    \noindent
    User U wants to save the completed trip. Once the trip is finished, the end-of-path screen appears, showing the distance traveled and the average speed, along with the option to submit a report on the path’s status and obstacles. Finally, U clicks the "Back to Dashboard" button and the trips is saved into Recent Trips.


    \vspace{3mm}
    \noindent
    \textbf{User inserts an already existing path}
    \vspace{1mm}
    
    \noindent
    User u enters the information for a new path, but the streets in the entered route are the same as those of an existing path. In this case, the application will notify the user with a warning that the entered streets correspond to those of an already existing path. The same thing happens whether the insertion is done manually or automatically

    \vspace{3mm}

    \noindent
    \textbf{Unregistered users} 
    \vspace{2mm}
    
    \noindent
    \textbf{User registers in the application}
    \vspace{1mm}
    
    \noindent
    User U, having navigated to the Login screen and clicks the "Registration" link presented on this page, is immediately redirected to the dedicated registration page. Here, U is required to complete the registration form by filling in all the necessary information: they enter a valid email address and define a secure password. After reviewing the entered data, U presses the final "Register" button. A verification email will then be sent to the provided email address, containing a verification code that has to be inserted in the application in the appropriate textbox, to confirm that the entered email address exists. Upon successful completion of the registration process, the system confirms the creation of the new account and redirects U back to the main Login screen, prompting them to perform their first access.

    \vspace{3mm}

    \noindent
    \textbf{All users} 
    \vspace{2mm}
  
    \noindent
    \textbf{User searches for a path} 
    \vspace{1mm}
    
    \noindent
     User U wants to find the best path to their destination. From the main map interface, U enters the desired origin and destination points into the floating search bar. The system calculates and displays the five highest-scoring paths on the map. Each path is shown with its associated score, including the visualization of the corresponding obstacles present along the path. The path with the highest score is highlighted in teal to indicate the recommended route, while the other alternatives are displayed for comparison, allowing U to choose the best option before starting the ride.

     \vspace{2mm}
  
    \noindent
    \textbf{User starts a navigation} 
    \vspace{1mm}
    
    \noindent
     After entering the origin and destination and selecting the desired path, user U wants to start navigation.
    From the main map interface, U clicks the "Start Navigation" button, and the navigation screen opens. On the map, a cursor representing U’s current position is displayed; at the top, the current speed and street name are shown, while at the bottom there is an "Exit Navigation" button to end the navigation.

\newpage
\subsubsection{Domain Class Diagram}
    \begin{figure}[!h]
        \centering
        \includegraphics[width=1.1\linewidth]{RASD/Images/Domain Class Diagram.png}
        \label{fig:placeholder}
    \end{figure}

    \noindent
    Notes:
    \begin{itemize}
    \item The \textit{creates} relationship indicates the publication of the Path, while \textit{reports} tracks the contributions regarding obstacles and status, which is then aggregated by the \textit{System} (Merging).
    \item The \textit{status} attribute in \textit{Path} is initially set by the user who published the path, and is later derived from the aggregation of the most recent \textit{PathReport}.
    \item The \textit{structures} attribute in \textit{Path} is a list of points in the map that allows the \textit{System} to draw the path. In manual mode they are manually inserted on the map by the user who publishes the path, while in automated mode they are automatically inserted by the \textit{System}.
    \end{itemize}

\newpage
\subsubsection{State Diagram}
    \noindent
    \textbf{Trip lifecycle} 
    \vspace{2mm}
    \begin{figure}[!h]
        \centering
        \includegraphics[width=1\linewidth]{RASD/Images/SD_trip_lifecycle.png}
        \label{fig:placeholder}
    \end{figure}
\noindent \\
The Trip lifecycle state diagram illustrates the two operational modes through which the entity is managed by the system.In the upper branch, dedicated to assisted navigation, the user enters the desired origin and destination; the system identifies a set of paths with the highest scores, from which the user selects their preferred route. Navigation is then initiated: during this phase, the system guides the user to the destination while simultaneously recording driving performance metrics. The flow concludes with a review phase, during which data can still be modified before completion.In parallel, the lower branch describes the automated mode, designed for passive tracking: here, the trip immediately enters the recording phase, allowing for intermediate pauses. The recording must then be stopped to proceed to the final validation of the acquired data.Both flows finally converge in the concluding saving state, ensuring that, regardless of the input mode, all metrics, coordinates, and reports are correctly persisted in the system database.
    \newpage
    \vspace{2mm}
    \noindent
    \textbf{Path lifecycle} 
    \vspace{2mm}
    \begin{figure}[!h]
        \centering
        \includegraphics[width=1\linewidth]{RASD/Images/SD_path_lifecycle.png}
        \label{fig:placeholder}
    \end{figure}

\noindent
The Path lifecycle state diagram illustrates the process of creating and publishing a new path within the system. Starting from the initial state, the flow branches into two distinct operational modes.The manual mode allows the user to compose the path by interactively selecting existing streets of interest on the map and manually placing obstacle markers, with the ability to freely switch between these activities.Alternatively, the automated mode acquires the track via real-time GPS recording, handling temporary interruptions through the pause state.Both flows converge in the validation phase; in this state, the user can review and eventually modify the collected information before final confirmation. The cycle concludes with the publication of the path, or with the discard of the draft if not confirmed.

\newpage
    \vspace{2mm}
    \noindent
        \textbf{User Profile lifecycle} 
    \vspace{2mm}
    \begin{figure}[!h]
        \centering
        \includegraphics[width=1\linewidth]{RASD/Images/SD_user_profile_lifecycle.png}
        \label{fig:placeholder}
    \end{figure}

\noindent
The lifecycle begins when a visitor without an account completes registration, creating a persistent profile in the system that starts in a signed-out state. Entering credentials activates the working session, allowing the user to modify personal data or undertake new trips to improve their performance statistics. Upon finishing these activities, logging out returns the user to an inactive state while keeping the updated data saved for future access. The cycle ends permanently only if the user chooses to delete their account, an action that irreversibly removes the profile from the system.


\subsection{Product funtions}
    \noindent
    This section outlines the principal functionalities of the platform.

    \vspace{2mm}
    \noindent
    \textbf{Sign up \& Login} 
    \vspace{1mm}
    
    \noindent
     To sign up for the system, a user must provide his name, an email address, and a password. The system checks whether there is an existing account with the same email address: if so, the user is notified that the account already exists; otherwise, the registration proceeds and a verification code is sent to the provided email address. To complete the registration, the user must enter the verification code in the designated textbox.
     To log in, a user must enter the email and the associated password, and the system checks if the credentials match with an existing account: if so, the log in is completed; otherwise, the user is notified that the credentials are wrong.

    \vspace{2mm}
    \noindent
    \textbf{Path search and visualization} 
    \vspace{1mm}
    
    \noindent
     The system allows users to find the best paths that connect two locations. They must enter the starting point and destination, and the system will display on the map the paths with the highest score, along with their related obstacles. It is also possible to access the information screen of the selected path.

    \vspace{2mm}
    \noindent
    \textbf{Publication of a path in manual mode} 
    \vspace{1mm}
    
    \noindent
     The system allows registered users to publish new paths by manually providing the required information.
    The user draws the path and adds obstacles directly on the map, and then enters the remaining information.

    \vspace{2mm}
    \noindent
    \textbf{Publication of a path in automated mode} 
    \vspace{1mm}
    
    \noindent
     The system allows registered users to publish new paths by providing the required information automatically.
    The system retrieves data from GPS and sensors (accelerometer, gyroscope) while the user is biking, and automatically fills in all the necessary information, and the user can choose to accept, modify, or discard it.

    \vspace{2mm}
    \noindent
    \textbf{Navigation} 
    \vspace{1mm}
    
    \noindent
     After the user has selected the desired path from the map and starts navigation, the system provides a real-time representation of the user’s position on the map and also displays the current speed and the street being traveled.

    \vspace{2mm}
    \noindent
    \textbf{Trip tracking and analysis} 
    \vspace{1mm}
    
    \noindent
    The system allows registered users to keep track of their trips. After completing the trip, a screen will appear showing the trip statistics, and once the confirmation button is clicked, the system saves the trip and, if available, enriches it with meteorological information (weather, wind, temperature) retrieved from an external service.

    \vspace{2mm}
    \noindent
    \textbf{Path reporting} 
    \vspace{1mm}
    
    \noindent
    The system allows registered users to submit a report regarding the status and the presence of obstacles on an existing path.
    After completing a trip associated with the path in question, the user can modify the status and add or remove any obstacles, and once the confirmation button is clicked, the report will be sent to the system.

    \vspace{2mm}
    \noindent
    \textbf{Data Aggregation and Merging} 
    \vspace{1mm}
    
    \noindent
    The system collects path reports from multiple users and merges conflicting information. It calculates the current status of a path based on the freshness of the reports and the number of confirming data obtained.

\subsection{User characteristics}

    \subsubsection{Registered users}
    \noindent
    Registered users are able to login to the platform. Their profiles contain personal information and some statistics such as total kilometers travelled, total trips completed, total paths published and the average speed of their trips. In addition to all the features available to unregistered users, they can publish new paths (in either manual or automated mode) and submit reports regarding existing paths. Furthermore, they can record their trips to populate their history and view detailed statistics at the end of the navigation.

    \subsubsection{Unregistered users}
    \noindent
    Unregistered users include casual users, tourists, or people who simply want to find a bike path without tracking their data. They can search for the best paths between two locations, select a desired path, and start the real-time navigation to follow the path, but without the ability to save the trip or view post-activity statistics.

\newpage
\subsection{Assumptions, dependencies and constraints}

    \subsubsection{Domain Assumptions}
    \noindent
    This section describes the fundamental assumptions regarding the environment and the behavior of all entities interacting with the system. These assumptions simplify the system’s design and implementation by setting clear expectations for the users, and the platform itself, ensuring the system functions efficiently within its intended contex

    \begin{description}
        \item[\textbf{[D1]}] The user’s device is equipped with a GPS, an accelerometer and a gyroscope.
        \item[\textbf{[D2]}] The user has a working internet connection.
        \item[\textbf{[D3]}] The user provides truthful information regarding the path to be published.
        \item[\textbf{[D4]}] The user provides truthful reports regarding paths.
        \item[\textbf{[D5]}] The external weather information service is reliable.
        \item[\textbf{[D6]}] The external map service is available and provides accurate topographical data.
    \end{description}

    \subsubsection{Dependencies}
    \noindent
    \begin{itemize}
    \item \textbf{Map Service}: it's responsible for providing the map that enables the graphical visualization of the paths. This is a critical dependency: if this service is unavailable, the core functionalities of the system cannot be performed.
    \item \textbf{Weather Service}: it's responsible for adding weather information to the users’ trips. This is a non-critical dependency: if this service is unavailable, the system will continue to function normally, simply omitting the weather information from the trip details. 
    \end{itemize}

%------------------------------------------------------------------------------------------------------------------------------------------------
\clearpage
{\color{Blue}{\section{Specific Requirements}}}
\label{sect:requirements}
\subsection{External Interface Requirements}
Here is shown the mockups of the UI. With the exception of the home page, which is shown from both the guest’s and the user’s perspectives, all other screens are displayed only from the user’s perspective, since the only differences are already shown on the home page.
\vspace{3mm}
\newline
\textbf{Home page} \newline
The \emph{Home page} acts as a dashboard for exploration. It assumes the cyclist immediately wants to either find a safe path to ride or start recording their own journey. On the left, the home page is shown from the user’s perspective; on the right, from the guest’s perspective

\begin{figure}[H]
    \centering
    \includegraphics[width=1\linewidth]{RASD/Images/Mockup/HomePage.png}
    \caption{Home page}
    \label{fig:placeholder}
\end{figure}

\newpage
\noindent
\textbf{Search path} \newline
From the \emph{Home page}, when the cyclist inserts a destination, the five paths that connect origin and destination with the highest score will appear on the map. If the cyclist wants to change the origin to a location that is not the current location, the “Start Navigation” button will disappear.

\begin{figure}[H]
    \centering
    \includegraphics[width=0.6\linewidth]{RASD/Images/Mockup/StartNavigation.png}
    \caption{Start navigation screen}
    \label{fig:placeholder}
\end{figure}

\newpage
\noindent
\textbf{Navigation} \newline
From the \emph{Search path screen}, when the cyclist clicks “Start Navigation”, the navigation begin and the \emph{Navigation sreen} appears. When the cyclist reaches the destination, the navigation ends and the \emph{Trip completed screen} appears.

\begin{figure}[h!]
    \centering
    \subfloat[Navigation screen]{\includegraphics[width=0.5\textwidth]{RASD/Images/Mockup/Navigation.png}}\hfill
    \subfloat[Trip completed screen]{\includegraphics[width=0.5\textwidth]{RASD/Images/Mockup/NavigationCompleted.png}}
\end{figure}

\newpage
\noindent
\textbf{Manual Mode} \newline
From the \emph{Home page}, when the user clicks the “Manual Mode” button, the \emph{Draw Path screen} appears. Here, the user can draw a path by clicking on the map to add points, which the Map Service then connects to create a path that follows the street layout. By clicking “Obstacles”, the \emph{Insert Obstacles screen} appears, allowing the user to place obstacles directly on the map. Finally, clicking “Next Step” opens the \emph{Review \& Publish screen}, where the user sets the status and publishes the path by clicking “Publish Path”.
\begin{figure}[h!]
    \centering
    \subfloat[Draw path screen]{\includegraphics[width=0.33\textwidth]{RASD/Images/Mockup/DrawPath.png}}\hfill
    \subfloat[Insert obstacles screen]{\includegraphics[width=0.33\textwidth]{RASD/Images/Mockup/InsertObstacles.png}}
    \subfloat[Review \& publish screen]{\includegraphics[width=0.33\textwidth]{RASD/Images/Mockup/Review&Publish.png}}
\end{figure}

\newpage
\noindent

\newpage
\subsection{Functional Requirements}

\textbf{Sign up \& Login}
\begin{description}
    \item[\textbf{[R1]}] The platform shall allow new users to register an account by providing personal information.
    \item[\textbf{[R2]}] The platform shall allow registered users to log in by entering their registered email and password.
\end{description}

\noindent
\textbf{Map Explorer}
\begin{description}
    \item[\textbf{[R3]}] The platform shall display an interactive map centered on the user's current location.
    \item[\textbf{[R4]}] The platform shall allow users to input an "Origin" and "Destination" of the desired path.
    \item[\textbf{[R5]}] The platform shall allow users to swap the origin and destination inputs.
    \item[\textbf{[R6]}] The platform shall calculate and display the top 5 paths options ranked by score.
    \item[\textbf{[R7]}] The platform shall visualize the geometry, quality, scores and obstacles of the top 5 paths on the map.
    \item[\textbf{[R8]}] The platform shall allow users to start the navigation mode for the currently selected path.
    \item[\textbf{[R9]}] When the navigation is started, the platform shall visualize the user's position in real-time.
    \item[\textbf{[R10]}] When the navigation is started, the platform shall display a dynamic navigation overlay providing the name of the street being traversed and the estimated speed.

    \item[\textbf{[R11]}] the platform shall allow users to interrupt the active navigation mode at any time.

    \item[\textbf{[R12]}]  When the navigation is completed, the platform shall present to registered users a summary screen containing statistics and the map of the path taken.

    \item[\textbf{[R13]}] When the navigation is completed, the platform shall allow registered users to modify the path status and the obstacles.
\end{description}

\noindent
\textbf{Active Path Recording}
\begin{description}
    \item[\textbf{[R14]}] The platform shall allow registered users to start an active recording session.
    \item[\textbf{[R15]}] The platform shall allow registered users to pause and resume the recording of the active trip.
    \item[\textbf{[R16]}] The platform shall utilize device sensors to automatically detect obstacles.
    \item[\textbf{[R17]}] The platform shall display real-time metrics during the trip, including duration, current speed, distance traveled, and weather conditions.

    \item[\textbf{[R18]}] The platform shall allow registered users to stop the recording and immediately transition to the summary screen to review the captured data.

    \item[\textbf{[R19]}] The platform shall allow registered users to modify status and obstacles in the captured data summary.

    \item[\textbf{[R20]}] The platform shall allow registered users to publish the captured path along with its data.
\end{description}

\newpage
\noindent
\textbf{Manual Path Creation}
\begin{description}
    \item[\textbf{[R21]}] The platform shall allow registered users to manually draw the path on the map.
    \item[\textbf{[R22]}] The platform shall allow registered users to manually place obstacle markers along the drawn path.
    \item[\textbf{[R23]}] The platform shall allow registered users to set the status of the drawn path.
    \item[\textbf{[R24]}] The platform shall allow registered users to publish the drawn path along with its data.
\end{description}

\noindent
\textbf{Trips \& Published Paths}
\begin{description}
    \item[\textbf{[R25]}] The platform shall display a dashboard listing the user's historical trips and their data, sorted chronologically.
    \item[\textbf{[R26]}] The platform shall allow registered users to view a list of paths published by them.
    \item[\textbf{[R27]}] The platform shall allow registered users to edit the information of any of their published paths.
    \item[\textbf{[R28]}] The platform shall allow registered users to delete any of their published paths.
\end{description}

\noindent
\textbf{User Profile}
\begin{description}
    \item[\textbf{[R29]}] The platform shall display a user profile screen summarizing the main aggregate statistics.
    \item[\textbf{[R30]}] The platform shall allow registered users to edit their profile information.
\end{description}



\newpage
\subsubsection{Use Cases Diagrams}
    \begin{figure}[H]
        \centering
        \includegraphics[width=1\linewidth]{RASD/Images/UseCaseDiagrams/GuestCyclistUser.jpg}
        \caption{General Utilities Use Case Diagram}
        \label{fig:placeholder}
    \end{figure}

This Use Case Diagram depicts the system’s functionalities organized through a generalization hierarchy. The abstract actor Cyclist represents the baseline capability of searching for paths, which is common to all users. Two actors inherit from this base:

\begin{itemize}
    \item The Guest (unauthenticated), who extends the base capability with access-control operations.
    \item The User (authenticated), who extends the base capability with the full suite of platform features.
\end{itemize}

Finally, the diagram illustrates the participation of external services (Map and Weather Services) in supporting these functionalities.
\newpage
\begin{figure}[H]
    \centering
    \includegraphics[width=1\linewidth]{RASD/Images/UseCaseDiagrams/ManualMode.jpg}
    \caption{Manual Mode Use Case Diagram}
    \label{fig:manual_mode}
\end{figure}

\noindent
This diagram illustrates the manual path creation process centered on the 'Publish Path' use case. This main activity mandatorily includes both 'Draw Path' (to define the route geometry) and 'Add Status' (to assign a rating). Additionally, the process can be optionally extended by 'Add Obstacle' to allow the user to report specific issues on the map.

\vspace{15mm}

\begin{figure}[H]
    \centering
    \includegraphics[width=1\linewidth]{RASD/Images/UseCaseDiagrams/AutomatedMode.jpg}
    \caption{Automated Mode Use Case Diagram}
    \label{fig:placeholder}
\end{figure}
\noindent
This diagram illustrates the automated path acquisition process. The Device Sensors actor refers to the device’s GPS, accelerometer, and gyroscope, which are essential for tracking the user's movement and detecting path quality in real-time.
\newpage

\subsubsection{Use Cases}

\setcounter{table}{0}
%tabella
\begin{table}[H]
    \centering
    \setlength\extrarowheight{2pt}
    \renewcommand{\arraystretch}{1.2}
    \begin{tabularx}{\textwidth}{|>{\bfseries}l|X|}
    \hline
    UC-01 & Sign up \\ \hline
    Actors & Guest\\ \hline
    Entry Condition &  The Guest is not currently logged in.\\ \hline
    Event Flow & 
    \begin{enumerate}
      \item The Guest opens the login screen.
      \item The Guest clicks "sign up".
      \item The Guest enters the required information.
      \item The Guest presses the "sign up" button.
      \item The platform validates the entered information.
      \item The platform sends a confirmation email to the provided email address.
      \item The platform shows a message to the guest to confirm the registration.
      \item The Guest confirms the registration by clicking on the link in the email. 
      \item The platform confirms the registration and activates the account.
      \item The Guest is redirected to the platform's homepage as a logged-in User.
    \end{enumerate}
    \\ \hline
    Exit Condition &  A new account is created and the Guest is now logged in as a User.\\ \hline
    Exception & 
    \begin{itemize}
        \item The Guest provides incorrect information.
        \item The Guest does not confirm the registration.
        \item The email is already in use.
    \end{itemize}
    \\ \hline
    \end{tabularx}
    \caption{Use Case UC-01}
    \label{UC-01}
\end{table}

\newpage
\begin{table}[H]
    \centering
    \setlength\extrarowheight{2pt}
    \renewcommand{\arraystretch}{1.2}
    \begin{tabularx}{\textwidth}{|>{\bfseries}l|X|}
    \hline
    UC-02 & Login \\ \hline
    Actors & Guest\\ \hline
    Entry Condition & The Guest is not currently logged in.\\ \hline
    Event Flow & 
    \begin{enumerate}
      \item The Guest opens the login screen.
      \item The Guest enters login credentials.
      \item The Guest clicks the "login" button.
      \item The platform verifies login credentials.
      \item The platform creates a login session for the guest.
      \item The Guest is redirected to the platform’s homepage as a logged-in User.
    \end{enumerate}
    \\ \hline
    Exit Condition & The Guest is successfully authenticated as a User.\\ \hline
    Exception & 
    \begin{itemize}
        \item The credentials provided are incorrect.
    \end{itemize}
    \\ \hline
   \end{tabularx}
   \caption{Use Case UC-02}
    \label{UC-02}
\end{table}

\begin{table}[H]
    \centering
    \setlength\extrarowheight{2pt}
    \renewcommand{\arraystretch}{1.2}
    \begin{tabularx}{\textwidth}{|>{\bfseries}l|X|}
    \hline
    UC-03 & Search Path \\ \hline
    Actors & Cyclist, Map Service\\ \hline
    Entry Condition & The Cyclist is viewing the platform’s homepage\\ \hline
    Event Flow & 
    \begin{enumerate}
      \item The platform loads the visual map from the Map Service.
      \item The Cyclist enters the origin.
      \item The Cyclist enters the destination.
      \item The platform retrieves the locations using the Map Service.
      \item The platform identifies the best 5 paths among those published.
      \item The platform displays the best 5 paths on the map.
      \item The Cyclist selects the desired path.
    \end{enumerate}
    \\ \hline
    Exit Condition &  The desired path is selected and the option to view its details is available to the Cyclist.\\ \hline
    Exception & 
    \begin{itemize}
        \item None.
    \end{itemize}
    \\ \hline
    \end{tabularx}
    \caption{Use Case UC-03}
    \label{UC-03}
\end{table} 

\begin{table}[H]
    \centering
    \setlength\extrarowheight{2pt}
    \renewcommand{\arraystretch}{1.2}
    \begin{tabularx}{\textwidth}{|>{\bfseries}l|X|}
    \hline
    UC-04 & Start Navigation \\ \hline
    Actors & Cyclist\\ \hline
    Entry Condition & The Cyclist has selected the desired path and the origin is the current location\\ \hline
    Event Flow & 
    \begin{enumerate}
      \item The platform shows the "Start Navigation" button.
      \item The Cyclist clicks "Start Navigation".
      \item The platform displays the selected path on the map.
      \item The platform guides the Cyclist in real-time.
      \item The platform detects that the Cyclist has reached the destination.
    \end{enumerate}
    \\ \hline
    Exit Condition &  The platform displays the trip summary.\\ \hline
    Exception & 
    \begin{itemize}
        \item None.
    \end{itemize}
    \\ \hline
   \end{tabularx}
   \caption{Use Case UC-04}
    \label{UC-04}
\end{table}

\begin{table}[H]
    \centering
    \setlength\extrarowheight{2pt}
    \renewcommand{\arraystretch}{1.2}
    \begin{tabularx}{\textwidth}{|>{\bfseries}l|X|}
    \hline
    UC-05 & Save trip \\ \hline
    Actors & User, Weather Service\\ \hline
    Entry Condition & The user is logged in and the navigation is completed.\\ \hline
    Event Flow & 
    \begin{enumerate}
      \item The platform retrieves weather data from the Weather Service.
      \item The platform displays the "Trip Completed" screen with statistics, weather info, and rating options. 
      \item The User selects the path status.
      \item The User clicks "Modify Obstacles" to report obstacles on the map.
      \item The User clicks "Save Trip".
      \item The platform saves the trip in the User's "Recent Trips" list.
      \item The platform automatically submits the path report to the system.
      \item The User is redirected to the "Recent Trips" screen.
    \end{enumerate}
    \\ \hline
    Exit Condition & The trip is saved and the "Recent Trips" screen is displayed to the User.\\ \hline
    Exception & 
    \begin{itemize}
        \item None.
    \end{itemize}
    \\ \hline
   \end{tabularx}
   \caption{Use Case UC-05}
    \label{UC-05}
\end{table}

\begin{table}[H]
    \centering
    \setlength\extrarowheight{2pt}
    \renewcommand{\arraystretch}{1.2}
    \begin{tabularx}{\textwidth}{|>{\bfseries}l|X|}
    \hline
    UC-06 & View Trips \\ \hline
    Actors & User\\ \hline
    Entry Condition & The User is logged in.\\ \hline
    Event Flow & 
    \begin{enumerate}
      \item The User opens the "Recent Trips" screen.
      \item The platform displays the list of saved trips.
      \item The User views the details of the trips in the list.
    \end{enumerate}
    \\ \hline
    Exit Condition &  The detailed view of the saved trips is displayed.\\ \hline
    Exception & 
    \begin{itemize}
        \item None.
    \end{itemize}
    \\ \hline
   \end{tabularx}
   \caption{Use Case UC-06}
    \label{UC-06}
\end{table}

\begin{table}[H]
    \centering
    \setlength\extrarowheight{2pt}
    \renewcommand{\arraystretch}{1.2}
    \begin{tabularx}{\textwidth}{|>{\bfseries}l|X|}
    \hline
    UC-07 & Edit Profile \\ \hline
    Actors & User\\ \hline
    Entry Condition & The User is logged in.\\ \hline
    Event Flow & 
    \begin{enumerate}
      \item The User opens the "Profile" screen.
      \item The User clicks on "Edit Profile".
      \item The User modifies the information.
      \item The User clicks "Confirm Changes"
      \item The platform validates the new data.
      \item The platform updates the profile and displays the success message.
    \end{enumerate}
    \\ \hline
    Exit Condition &  The User profile is updated with the new information.\\ \hline
    Exception & 
    \begin{itemize}
        \item The User leaves a mandatory field empty and clicks "Confirm"
    \end{itemize}
    \\ \hline
   \end{tabularx}
   \caption{Use Case UC-07}
    \label{UC-07}
\end{table}

\begin{table}[H]
    \centering
    \setlength\extrarowheight{2pt}
    \renewcommand{\arraystretch}{1.2}
    \begin{tabularx}{\textwidth}{|>{\bfseries}l|X|}
    \hline
    UC-08 & Delete path \\ \hline
    Actors & User\\ \hline
    Entry Condition & The User is logged in.\\ \hline
    Event Flow & 
    \begin{enumerate}
      \item The User opens the "Published path" screen.
      \item The User selects the path to delete.
      \item The platform displays the path details.
      \item The User clicks the "Delete" button.
      \item The platform asks for confirmation.
      \item The User confirms the deletion.
      \item The platform removes the path and updates the list.
    \end{enumerate}
    \\ \hline
    Exit Condition &  The path is permanently deleted from the published paths list.\\ \hline
    Exception & 
    \begin{itemize}
        \item None.
    \end{itemize}
    \\ \hline
   \end{tabularx}
   \caption{Use Case UC-08}
    \label{UC-08}
\end{table}

\begin{table}[H]
    \centering
    \setlength\extrarowheight{2pt}
    \renewcommand{\arraystretch}{1.2}
    \begin{tabularx}{\textwidth}{|>{\bfseries}l|X|}
    \hline
    UC-09 & Publish Path (Manual Mode) \\ \hline
    Actors & User\\ \hline
    Entry Condition & The User is logged in and is viewing the homepage.\\ \hline
    Event Flow & 
    \begin{enumerate}
      \item The User clicks the "+" button to create a new path.
      \item The platform loads the "Draw Mode" interface.
      \item The User draw the path.
      \item The User clicks the "Next Step" button.
      \item The platform displays the "Review \& Publish" screen.
      \item The User selects the status.
      \item The User clicks "Publish Path".
      \item The platform validates and makes the path available to the community.
    \end{enumerate}
    \\ \hline
    Exit Condition &  The path is published and appears in the User's "Published Paths" screen.\\ \hline
    Exception &
    \begin{itemize}
        \item  The platform detects that an identical path already exists.
    \end{itemize}
    \\ \hline
   \end{tabularx}
   \caption{Use Case UC-09}
    \label{UC-09}
\end{table}

\begin{table}[H]
    \centering
    \setlength\extrarowheight{2pt}
    \renewcommand{\arraystretch}{1.2}
    \begin{tabularx}{\textwidth}{|>{\bfseries}l|X|}
    \hline
    UC-10 & Draw Path \\ \hline
    Actors & User, Map Service\\ \hline
    Entry Condition & The User is in the 'Draw Path' mode within the manual creation interface and the interactive map is loaded.\\ \hline
    Event Flow & 
    \begin{enumerate}
      \item The User taps on the map to place the starting point.
      \item The User continues tapping to trace the path segments.
      \item The Map Service renders the route line in real-time.
      \item The platform calculates the "Total Km" dynamically.
      \item The User completes the tracing.
    \end{enumerate}
    \\ \hline
    Exit Condition &  The geometric path is defined on the map.\\ \hline
    Exception & 
    \begin{itemize}
        \item The platform fails to load the map tiles from the Map Service.
    \end{itemize}
    \\ \hline
   \end{tabularx}
   \caption{Use Case UC-10}
    \label{UC-10}
\end{table}

\begin{table}[H]
    \centering
    \setlength\extrarowheight{2pt}
    \renewcommand{\arraystretch}{1.2}
    \begin{tabularx}{\textwidth}{|>{\bfseries}l|X|}
    \hline
    UC-11 & Add Obstacle \\ \hline
    Actors & User, Map Service\\ \hline
    Entry Condition & The User is in the manual creation interface, the interactive map is loaded, and the path is partially or fully drawn.\\ \hline
    Event Flow & 
    \begin{enumerate}
      \item The User clicks on the "Obstacles" tab.
      \item The platform switches to "Obstacle Mode".
      \item The User taps a specific point along the drawn path.
      \item The platform places an obstacle marker on the coordinate using the Map Service.
    \end{enumerate}
    \\ \hline
    Exit Condition & An obstacle is flagged on the map.\\ \hline
    Exception & 
    \begin{itemize}
        \item The user adds an obstacle outside the path
    \end{itemize}
    \\ \hline
   \end{tabularx}
   \caption{Use Case UC-11}
    \label{UC-11}
\end{table}

\begin{table}[H]
    \centering
    \setlength\extrarowheight{2pt}
    \renewcommand{\arraystretch}{1.2}
    \begin{tabularx}{\textwidth}{|>{\bfseries}l|X|}
    \hline
    UC-12& Start Automated Path Acquisition \\ \hline
    Actors & User, Device Sensors\\ \hline
    Entry Condition & The User is logged in and is viewing the homepage\\ \hline
    Event Flow & 
    \begin{enumerate}
      \item The User clicks the "Lightning" button to start the acquisition.
      \item The platform initializes the "Automated Mode" interface.
      \item The platform loads the User's current location on the map using the gps from the Device Sensors.
      \item The platform starts the timer.
      \item The platform begins recording distance, speed and obstacles in real-time using the Device Sensors.
      \item The platform continuously updates the metrics (distance, obstacles, speed).
    \end{enumerate}
    \\ \hline
    Exit Condition & The path recording is active and the "Automated Mode" interface is displayed.\\ \hline
    Exception & 
    \begin{itemize}
        \item None
    \end{itemize}
    \\ \hline
   \end{tabularx}
   \caption{Use Case UC-12}
    \label{UC-12}
\end{table}

\begin{table}[H]
    \centering
    \setlength\extrarowheight{2pt}
    \renewcommand{\arraystretch}{1.2}
    \begin{tabularx}{\textwidth}{|>{\bfseries}l|X|}
    \hline
    UC-13& Pause Path Acquisition \\ \hline
    Actors & User\\ \hline
    Entry Condition & The User is logged in and the "Automated Mode" is active\\ \hline
    Event Flow & 
    \begin{enumerate}
      \item The User clicks the "Pause" button.
      \item The platform suspends the timer.
      \item The platform stops recording coordinates and metrics.
      \item The platform processes the information.
      \item The platform updates the interface to show the "Paused" state.
    \end{enumerate}
    \\ \hline
    Exit Condition & The "Automated Mode" is temporarily suspended.\\ \hline
    Exception & 
    \begin{itemize}
        \item None
    \end{itemize}
    \\ \hline
   \end{tabularx}
   \caption{Use Case UC-13}
    \label{UC-13}
\end{table}

\begin{table}[H]
    \centering
    \setlength\extrarowheight{2pt}
    \renewcommand{\arraystretch}{1.2}
    \begin{tabularx}{\textwidth}{|>{\bfseries}l|X|}
    \hline
    UC-14& Stop Path Acquisition \\ \hline
    Actors & User, Map Service\\ \hline
    Entry Condition & The User is logged in and the "Automated Mode" is active (or paused)\\ \hline
    Event Flow & 
    \begin{enumerate}
      \item The User clicks the "Stop" button.
      \item The platform stops the timer.
      \item The platform stops the data recording.
      \item The platform processes the recorded data to calculate final metrics.
      \item The platform extracts the street names using the Map Service.
      \item The platform displays the "Trip Summary" interface.
    \end{enumerate}
    \\ \hline
    Exit Condition &  The "Automated Mode" is stopped and the User is viewing the "Trip Summary"\\ \hline
    Exception & 
    \begin{itemize}
        \item None
    \end{itemize}
    \\ \hline
   \end{tabularx}
   \caption{Use Case UC-14}
    \label{UC-14}
\end{table}

\begin{table}[H]
    \centering
    \setlength\extrarowheight{2pt}
    \renewcommand{\arraystretch}{1.2}
    \begin{tabularx}{\textwidth}{|>{\bfseries}l|X|}
    \hline
    UC-15& Publish Path (Automated Mode) \\ \hline
    Actors & User\\ \hline
    Entry Condition & The User is viewing the "Trip Summary" interface with the recorded path data.\\ \hline
    Event Flow & 
    \begin{enumerate}
      \item The User clicks the "Publish Path" button.
      \item The Platform stores the path data (geometry, status, obstacles) in the system.
      \item The platform validates and makes the path available to the community..
    \end{enumerate}
    \\ \hline
    Exit Condition &  The path is published and appears in the User’s "Published Paths" screen.\\ \hline
    Exception & 
    \begin{itemize}
        \item The platform detects that an identical path already exists.
    \end{itemize}
    \\ \hline
   \end{tabularx}
   \caption{Use Case UC-15}
    \label{UC-15}
\end{table}


\newpage
\subsubsection{Sequence Diagrams}
The functions used in the sequence diagrams follow these conventional meanings:
\begin{itemize}
    \item \textbf{press(Button)}: Represents the action of a user pressing Button.
    \item \textbf{submit(Content)}: Indicates a user uploading Content to the platform.
    \item \textbf{showPage(Page)}: Displays the content of a specific Page.
    \item \textbf{showError(Error)}: Displays an error message indicating an issue that occurred.
    \item \textbf{check()}: Verifies the validity of the information provided by the user.
    \item \textbf{update(Data)}: Refers to internal processes that update database records with Data.
    \item \textbf{OpenApplication()}: Rappresent the action of a user opens the application.
    \item \textbf{getMapData()}: Represents the system requesting the visual map components from the external service.
    \item \textbf{mapTiles(Tiles)}: Delivers the graphical data required to render the map interface on the screen.
    \item \textbf{enter(Location)}: Represents the action of a user inputting a starting point or destination address.
    \item \textbf{findBestPaths(Origin, Destination)}: Identifies and returns the 5 paths with the best performance scores connecting the selected locations.
    \item \textbf{selectPath(Path)}: Represents the user choosing one of the displayed paths to view its details or start navigation.
    \item \textbf{startTracking()}: Commands the device sensors (GPS) to begin sending real-time location updates.
    \item \textbf{stopTracking()}: Commands the device sensors to cease location updates
    \item \textbf{getWeather(CurrentLocation)}: Requests meteorological data (temperature, conditions) for the specific coordinates where the trip ended.
    \item \textbf{wetherData()}: Returns the requested weather information from the external service.
    \item \textbf{insert(Object)}: Represents the user placing a marker on the map.
    \item \textbf{initializeSensors()}: Internal process that activates and configures the device's hardware to start receiving real-time data.
    \item \textbf{monitorRoadConditions()}: Analyzes data from the accelerometer and gyroscope to automatically detect potholes during the ride.
    \item \textbf{displayInstruction()}: Updates the user interface with visual navigation cues and the current position on the map.
    \item \textbf{fetchData(Criteria)}: A generic system request to retrieve specific datasets from the database required to populate or plot elements on the user interface.
    \item \textbf{showMessage(Argument)}: Shows a general message on the screen.
    \item \textbf{startTracking()}: Activates the continuous recording and initiates the loop.
    \item \textbf{monitorRoadConditions()}:Processes live sensor data to understand the path environment.
    \item \textbf{stopTracking()}: Commands the device sensors to cease location updates.

    
\end{itemize}
\newpage
\vspace{2mm}
\noindent
    \textbf{[SD1] Sign up Sequence Diagram} 
\vspace{2mm}
\begin{figure}[!h]
    \centering
    \includegraphics[width=0.8\linewidth]{RASD/Images/SEQD_SignUp.png}
    \label{fig:placeholder}
\end{figure}
\newpage
\vspace{2mm}
\noindent
    \textbf{[SD2] Login Sequence Diagram} 
\vspace{2mm}
\begin{figure}[!h]
    \centering
    \includegraphics[width=0.8\linewidth]{RASD/Images/SEQD_Login.png}
    \label{fig:placeholder}
\end{figure}
\newpage
\vspace{2mm}
\noindent
    \textbf{[SD3] Search Path Sequence Diagram} 
\vspace{2mm}
\begin{figure}[!h]
    \centering
    \includegraphics[width=0.8\linewidth]{RASD/Images/SEQD_search_path.png}
    \label{fig:placeholder}
\end{figure}
\newpage
\vspace{2mm}
\noindent
    \textbf{[SD4] Start Navigation Sequence Diagram} 
\vspace{2mm}
\begin{figure}[!h]
    \centering
    \includegraphics[width=0.8\linewidth]{RASD/Images/SEQD_start_navigation.png}
    \label{fig:placeholder}
\end{figure}
\vspace{2mm}
\newpage
\noindent
    \textbf{[SD5] Save Trip Sequence Diagram} 
\vspace{2mm}
\begin{figure}[!h]
    \centering
    \includegraphics[width=0.8\linewidth]{RASD/Images/SEQD_save_trip.png}
    \label{fig:placeholder}
\end{figure}
\newpage
\noindent
    \textbf{[SD6] Trip View Sequence Diagram} 
\vspace{2mm}
\begin{figure}[!h]
    \centering
    \includegraphics[width=0.8\linewidth]{RASD/Images/SEQD_view_trip.png}
    \label{fig:placeholder}
\end{figure}
\newpage
\noindent
    \textbf{[SD7] Edit Profile Sequence Diagram} 
\vspace{2mm}
\begin{figure}[!h]
    \centering
    \includegraphics[width=0.8\linewidth]{RASD/Images/SEQD_edit_profile.png}
    \label{fig:placeholder}
\end{figure}
\newpage
\noindent
    \textbf{[SD8] Delete Path Sequence Diagram} 
\vspace{2mm}
\begin{figure}[!h]
    \centering
    \includegraphics[width=0.8\linewidth]{RASD/Images/SEQD_delete_path.png}
    \label{fig:placeholder}
\end{figure}
\newpage
\noindent
    \textbf{[SD9] Publish Path Sequence Diagram} 
\vspace{2mm}
\begin{figure}[!h]
    \centering
    \includegraphics[width=0.8\linewidth]{RASD/Images/SEQD_publish_path_manual_mode.png}
    \label{fig:placeholder}
\end{figure}
\newpage
\noindent
    \textbf{[SD10] Draw Path Sequence Diagram} 
\vspace{2mm}
\begin{figure}[!h]
    \centering
    \includegraphics[width=0.8\linewidth]{RASD/Images/SEQD_draw_path.png}
    \label{fig:placeholder}
\end{figure}
\newpage
\noindent
    \textbf{[SD11] Add Obstacle Sequence Diagram} 
\vspace{2mm}
\begin{figure}[!h]
    \centering
    \includegraphics[width=0.8\linewidth]{RASD/Images/SEQD_add_obstacle.png}
    \label{fig:placeholder}
\end{figure}
\newpage
\noindent
    \textbf{[SD12] Start Automated Path Acquisition Sequence Diagram} 
\vspace{2mm}
\begin{figure}[!h]
    \centering
    \includegraphics[width=0.8\linewidth]{RASD/Images/SEQD_start_automated_path_acquisition.png}
    \label{fig:placeholder}
\end{figure}
\newpage
\noindent
    \textbf{[SD13] Pause Path Acquisition Sequence Diagram} 
\vspace{2mm}
\begin{figure}[!h]
    \centering
    \includegraphics[width=0.8\linewidth]{RASD/Images/SEQD_pause_path_acquisition.png}
    \label{fig:placeholder}
\end{figure}
\newpage
\noindent
    \textbf{[SD14] Stop Path Acquisition Sequence Diagram} 
\vspace{2mm}
\begin{figure}[!h]
    \centering
    \includegraphics[width=0.8\linewidth]{RASD/Images/SEQD_stop_path_acquisition.png}
    \label{fig:placeholder}
\end{figure}
\newpage
\noindent
    \textbf{[SD15] Publish Path Automatated Mode Sequence Diagram} 
\vspace{2mm}
\begin{figure}[!h]
    \centering
    \includegraphics[width=0.8\linewidth]{RASD/Images/SEQD_publish_path_automated_mode.png}
    \label{fig:placeholder}
\end{figure}

\newpage
\subsection{Requirements Mapping}
\begin{table}[H]
    \centering
    % Aumenta lo spazio tra le righe per leggibilità
    \renewcommand{\arraystretch}{1.5}
    
    % Definizione tabella: larghezza totale testo, bordo verticale, colonna stretta per ID, colonna larga X per testo
    \begin{tabularx}{\textwidth}{|l X|}
        \hline
        % Riga Intestazione (G1) che occupa due colonne
        \multicolumn{2}{|l|}{\textbf{G1}: Registered users should have access to performance metrics about their trips.} \\
        \hline
        & \\[-1em] % Piccolo aggiustamento spaziatura
        
        % Domain Assumptions
        {[D12]} & The user’s device is equipped with a GPS, an accelerometer and a gyroscope. \\
        {[D2]} & The user has a working internet connection. \\
        {[D6]} & The external map service is available and provides accurate topographical data.\\
        & \\
        \hline
        & \\[-1em]
        
        % Requirements
        {[R1]} & The platform shall allow new users to register an account by providing personal information. \\
        {[R2]} & The platform shall allow registered users to log in by entering their registered email and password. \\
        {[R8]} & The platform shall allow users to start the navigation mode for the currently selected path. \\
        {[R12]} & When the navigation is completed, the platform shall present to registered users a summary screen
        containing statistics and the map of the path taken. \\
        {[R17]} & The platform shall display real-time metrics during the trip, including duration, current speed,
        distance traveled, and weather conditions. \\
        {[R18]} & The platform shall allow registered users to stop the recording and immediately transition to the summary screen to review the captured data. \\
        {[R25]} & The platform shall display a dashboard listing the user’s historical trips and their data, sorted chronologically.\\
        {[R29]} & The platform shall display a user profile screen summarizing the main aggregate statistics. \\
        & \\
        \hline
    \end{tabularx}
\end{table}

\newpage
\begin{table}[H]
    \centering
    % Aumenta lo spazio tra le righe per leggibilità
    \renewcommand{\arraystretch}{1.5}
    
    % Definizione tabella: larghezza totale testo, bordo verticale, colonna stretta per ID, colonna larga X per testo
    \begin{tabularx}{\textwidth}{|l X|}
        \hline
        % --- CORREZIONE QUI SOTTO ---
        % Uso 'p' con una larghezza calcolata (larghezza totale - margini interni) invece di 'l'
        % Questo obbliga il testo lungo ad andare a capo.
        \multicolumn{2}{|p{\dimexpr\textwidth-2\tabcolsep}|}{\textbf{G2}: When available, users would like to be informed about meteorological information related to the paths of interest.} \\
        \hline
        & \\[-1em] % Piccolo aggiustamento spaziatura
        
        % Domain Assumptions
        {[D2]} & The user has a working internet connection. \\
        {[D5]} & The external weather information service is reliable.\\
        & \\
        \hline
        & \\[-1em]
        
        % Requirements
        {[R1]} & The platform shall allow new users to register an account by providing personal information. \\
        {[R2]} & The platform shall allow registered users to log in by entering their registered email and password. \\
        {[R17]} & The platform shall display real-time metrics during the trip, including duration, current speed, distance traveled, and weather conditions. \\
        {[R25]} & The platform shall display a dashboard listing the user’s historical trips and their data, sorted chronologically. \\ % Corretto: rimosso '&' extra
        & \\
        \hline
    \end{tabularx}
\end{table}

%------------------------------------------------------------------------------------------------------------------------------------------------


\clearpage
{\color{Blue}{\section{Formal Analysis Using Alloy}}}
\label{sect:alloy}
\noindent
The objective of this section is to validate the critical aspects of the Best Bike Paths (BBP) system using Alloy, a lightweight formal method based on relational logic. In particular, the modeling effort concentrates on the robustness of the report submission mechanism, a core feature for maintaining data reliability within the platform.
\subsection{Signatures}
The signatures defined below represent the vocabulary used to describe the BBP ecosystem. Here, we outline the primary atoms of the model, including User, Trip, and Path, while also defining the dynamic elements such as Obstacle and PathStatus that are essential for tracking the system's evolution.

\lstinputlisting[language=alloy]{RASD/Alloy/Signatures.txt}


\newpage
\subsection{Facts}
The facts represent the constraints that the system must respect and are always true in the model. For the BBP application, we defined numerous facts to enforce data consistency, such as the necessity for a PathReport to be strictly linked to a Trip performed by the author. Furthermore, they dictate the logic for the dynamic evolution of the system, determining how PathStatus and Obstacles are updated based on the aggregation of user reports.
\lstinputlisting[language=alloy]{RASD/Alloy/Facts.txt}

\subsection{Assertions}
The assertions are the safety properties that we want to verify in the model. These conditions must hold true to ensure the platform avoids invalid or illogical states. For the BBP scenario, we verified several critical aspects, such as the strict requirement that a User must have actually completed a Trip on a specific Path before submitting a report (AuthorWasThere), as well as the constraint that limits users to a single report submission per time step to prevent spam.
\lstinputlisting[language=alloy]{RASD/Alloy/Assertions.txt}

\subsection{Model Visualization}
The following model was obtained with the following predicate:
\lstinputlisting[language=alloy]{RASD/Alloy/Predicate.txt}

\vspace{5mm}
\noindent
\textbf{Timeframe 0 (focus on discard trip):}

\begin{figure}[!h]
    \centering
    \includegraphics[width=1.1\linewidth]{RASD/Images/Alloy/Timeframe0.png}
    \label{fig:placeholder}
\end{figure}
\noindent
User3 takes a trip on Path1 but decides to discard it; therefore, no report is submitted. The status of all paths remains unchanged, and the set of obstacles remains the same.

\newpage
\textbf{Timeframe 1 (focus on tie and add/remove obstacles):}
\begin{figure}[!h]
    \centering
    \includegraphics[width=1.1\linewidth]{RASD/Images/Alloy/Timeframe1.png}
    \label{fig:placeholder}
\end{figure}

\noindent
User0 emits a report on Path0 with status = 'Manteinance Required'. \\
User1 emits a report on Path0 with status = 'Sufficient' and one obstacle (Obstacle3).\\
User2 emits a report on Path2 with status = 'Manteinance Required' and two obstacles (Obstacle1 and Obstacle4).\\
User3 emits a report on Path1 with status = 'Medium' and one obstacle (Obstacle 2), and it removes Obstacle0.\\
As a result, Path1 and Path2 take the status of the single report associated, while since there is a tie in the count of report statuses for Path0, it is assigned the lower of the two statuses (Maintenance Required).\\
In addition, Obstacle0 will be replaced with Obstacle3 in Path1, Obstacle1 and Obstacle4 have been added to Path2, and Obstacle2 has been added to Path1.

\newpage
\textbf{Timeframe 2 (focus on merging):}
\begin{figure}[!h]
    \centering
    \includegraphics[width=1.1\linewidth]{RASD/Images/Alloy/Timeframe2.png}
    \label{fig:placeholder}
\end{figure}
\\
User0 emits a report on Path2 with status = 'Medium'. \\
User1 emits a report on Path0 with status = 'Medium'.\\
User3 emits a report on Path2 with status = 'Sufficient'.\\
User2 makes a trip on Path2 but decides to discard it, therefore no report is submitted.\\
As a result, the status of Path2 is updated to Medium according to the majority vote, while the remaining paths are unaffected.

%------------------------------------------------------------------------------------------------------------------------------------------------
\clearpage
{\color{Blue}{\section{Use of Generative AI}}}
\label{sect:generativeai}
\subsection{AI-Based Support Tools}

While tools such as ChatGPT were utilized transversally to support drafting, code refinement, and logical requirement analysis, the specific phase of User Interface Prototyping relied on models specialized in image generation.

\noindent
Specifically, Google AI Studio was employed to interact with the Gemini model. Thanks to its multimodal capabilities, it was possible to translate the textual descriptions of Use Cases into high-fidelity visual mockups. This approach allowed for a rapid visualization of the spatial organization of elements (layout).

\subsection{Example of Interaction: Mockup Generation}

 A significant example of this interaction is represented by the definition of the Homepage. Initially, a prompt was provided to the model aimed at establishing the general structure and color palette, inspired by leading market competitors.

\begin{PromptBox}
\quad Generate a mockup for the homepage of a mobile GPS navigation application. The interface must be in "Dark Mode", featuring a dark night-time street map background similar to Google Maps. At the bottom of the screen, include a navigation bar with four sections: "Home", "Published", "Recent Trip", and "Profile". Position two round floating buttons stacked vertically immediately above this bar in the bottom right corner: one for "Manual Mode" and the other for "Automated Mode".

\vspace{0.3cm}

\end{PromptBox}

\begin{GeminiBad}
\quad Here is a proposal for the Homepage design in night mode. The generated image presents: \begin{itemize} \item A dark-toned cartographic background (midnight blue and charcoal) with highlighted roads. \item A bottom navigation bar featuring the four requested sections. \item The two action buttons ("Manual" and "Automated") positioned at the bottom right for easy thumb access. \end{itemize}

\textbf{} \\
    \begin{center}
    \includegraphics[height=6cm, keepaspectratio]{RASD/Images/WrongImage.png}
    \end{center}
    
\vspace{0.3cm}
\end{GeminiBad}

\noindent
Although the structure was correct, the first iteration lacked crucial elements for immediate usability, such as a clear indication of the user's position and the destination input field. A second refinement prompt was therefore necessary to integrate these details without altering the already approved visual style.

\vspace{0.3cm}
\begin{PromptBox}
\quad Based on the image just generated, modify only the top area and the center of the map. Add a floating search bar at the top of the screen specifically for entering a destination (with a placeholder like "Where to?"). Additionally, place a prominent cursor, such as a glowing blue arrow, in the center of the map to indicate the user's current position. Ensure that the rest of the existing interface, including the bottom navigation bar, remains unchanged.

\vspace{0.3cm}
\end{PromptBox}


\begin{GeminiGood}
\quad Updated the mockup by integrating the requested elements: \begin{itemize} \item \textbf{Search Bar:} Inserted at the top with rounded corners and a search icon, maintaining the transparency typical of modern UIs. \item \textbf{Position Cursor:} Added a blue indicator with a slight glow effect at the center of the map to ensure immediate contrast against the dark background. \end{itemize} The rest of the layout (navbar and mode buttons) remained unchanged.

\textbf{} \\
    \begin{center}
    \includegraphics[height=6cm, keepaspectratio]{RASD/Images/TrueImage.png}
    \end{center}
    
\vspace{0.3cm}

\end{GeminiGood}

\vspace{0.5cm}



\noindent
This two-phase process allowed for the separation of the structural layout definition (navbar, map, buttons) from the functional details (search bar, cursor). The model demonstrated the ability to maintain stylistic consistency between generations, integrating the requested modifications naturally within the pre-existing visual context.

\newpage
\subsection{Example of Interaction: Rules Identification}

The following interaction describes a consultation aimed at defining the Standard Compliance section. We engaged the AI to clarify the complex requirements related to the handling of continuous geolocation history and the integration of external map services. Through two prompts, we progressed from identifying the applicable external regulations (GDPR, Terms of Service) to defining the internal technical solutions required to ensure compliance without compromising the platform’s functionality.

\vspace{0.5cm}

\begin{PromptBox}
\quad I am developing a mobile app for cyclists that features real-time navigation and trip tracking, relying on third-party map services. Considering I will be handling continuous geolocation data, what are the primary privacy regulations and map licensing constraints I need to strictly follow?

\vspace{0.3cm}

\end{PromptBox}

\begin{ChatGPTBad}
\quad Based on your description, the platform operates at the intersection of sensitive user data and proprietary services. Here are the mandatory constraints:
\begin{enumerate}
    \item Privacy and Sensitive Data: Since the system handles Personally Identifiable Information (PII) and continuous geolocation history, this constitutes a significant privacy risk as it reveals user habits and private locations. You must strictly comply with the General Data Protection Regulation (GDPR) and potentially the California Consumer Privacy Act (CCPA).
    \item Map Provider Terms of Service (ToS):The dependency on external providers (e.g., OpenStreetMap or Google Maps) mandates rigorous adherence to their ToS. You are required to:
    \begin{itemize}
        \item Dynamically display proper copyright attributions.
        \item Strictly enforce technical limitations on offline data caching to avoid licensing violations.
    \end{itemize}
\end{enumerate}
    
\vspace{0.3cm}
\end{ChatGPTBad}



\vspace{0.5cm}

\noindent
The analysis of the regulatory landscape revealed a critical tension: the application requires continuous data collection for functionality, yet strictly limits data retention for compliance. Recognizing that simple adherence to rules is insufficient, we directed the next query toward defining the architectural drivers necessary to resolve this conflict at the structural level.

\vspace{0.3cm}
\newpage
\begin{PromptBox}
\quad Based on these regulatory constraints, how should I design the software architecture to ensure compliance? Specifically, I need to understand how to embed "Privacy-by-Design" principles into the core system and how to technically enforce the caching limitations, while keeping the architecture flexible enough to adapt to future changes.

\vspace{0.3cm}
\end{PromptBox}


\begin{ChatGPTGood}
\quad To ensure compliance without compromising functionality, the system must be architected with specific drivers:
\begin{itemize} \item \textbf{Privacy-by-Design Core:} You must embed privacy-by-design principles directly into the software core, rather than treating compliance as an external add-on. This minimizes the risk of exposing user habits. \item \textbf{Adaptability and Flexibility:} The system must be architected with the flexibility to adapt to this dual set of regulatory and licensing constraints. Regulations and ToS change frequently; a modular design will allow you to update compliance logic without rewriting the entire application.\item \textbf{Technical Enforcement:} The architecture must include a logic layer that automatically manages data retention, ensuring you strictly enforce technical limitations on offline data caching as required by the map providers.\end{itemize} The rest of the layout (navbar and mode buttons) remained unchanged.
    
\vspace{0.3cm}

\end{ChatGPTGood}


\vspace{0.5cm}

\noindent \textbf{Verification} \\
We verified the information generated by the AI by checking official online sources. The details regarding privacy regulations and map service policies were confirmed to be accurate and consistent with real-world requirements. Therefore, the content is reliable and has been validated for use in this project.

\subsection{Conclusion}
The integration of AI-based tools has significantly improved efficiency in the preparation of complex technical documents. Nevertheless, despite their advantages, it remains essential to adopt a critical approach and carefully review the produced content to ensure accuracy, coherence, and completeness. Achieving high-quality results often requires multiple iterations and precise refinements, guiding the process step by step. AI tools have proven to be valuable support instruments, but they should be used with careful supervision throughout all phases of document development.

%------------------------------------------------------------------------------------------------------------------------------------------------

\clearpage
{\color{Blue}{\section{Effort Spent}}}
\label{sect:effort}
\subsection*{Andrea Ferritti}

\begin{table}[H]
    \centering
\begin{tabular}{|l|c|}
        \hline
        \textbf{Section} & \textbf{Hours} \\ \hline
        Introduction & 10 \\ \hline
        Overall Description & 9\\ \hline
        Specific Requirements & 16\\ \hline
        Formal Analysis Using Alloy & 11 \\ \hline
        Use of Generative AI & 1 \\ \hline
    \end{tabular}
    \caption{Effort - Andrea Ferritti}
    \label{tab:effortferritti}
\end{table}

\subsection*{Jieke Lin}

\begin{table}[H]
    \centering
\begin{tabular}{|l|c|}
        \hline
        \textbf{Section} & \textbf{Hours} \\ \hline
        Introduction & 10 \\ \hline
        Overall Description & 8 \\ \hline
        Specific Requirements & 16\\ \hline
        Formal Analysis Using Alloy & 12 \\ \hline
        Use of Generative AI & 1 \\ \hline
    \end{tabular}
    \caption{Effort - Jieke Lin}
    \label{tab:effortlin}
\end{table}

\subsection*{Tommaso Tosca}

\begin{table}[H]
    \centering
\begin{tabular}{|l|c|}
        \hline
        \textbf{Section} & \textbf{Hours} \\ \hline
        Introduction & 9 \\ \hline
        Overall Description & 11\\ \hline
        Specific Requirements & 16\\ \hline
        Formal Analysis Using Alloy & 14 \\ \hline
        Use of Generative AI & 1 \\ \hline
    \end{tabular}
    \caption{Effort - Tommaso Tosca}
    \label{tab:efforttosca}
\end{table}


%------------------------------------------------------------------------------------------------------------------------------------------------




%\clearpage
%\addcontentsline{toc}{section}{References}
%\bibliographystyle{plain}
%\bibliography{main}
%------------------------------------------------------------------------------------------------------------------------------------------------




\end{document}
}
\else
    \let\commandsFileStandaloneBootstrap\relax
\fi
\commandsFileStandaloneBootstrap

\fi

\usepackage{tikz}

\usepackage{tabularx}
\usepackage[dvipsnames]{xcolor}
\usepackage{listings}
\usepackage{Alloy/alloy-style}


\usepackage[most]{tcolorbox}

\tcbset{
  chatstyle/.style={
    breakable,
    enhanced,
    rounded corners,    % angoli arrotondati
    arc=2mm,            % raggio arrotondamento
    boxrule=1pt,
    coltitle=black,      % Colore del testo del titolo
    fonttitle=\bfseries, % Titolo in grassetto
    attach title to upper,
    bottomtitle=1mm,
    toptitle=1mm,
    titlerule=0.5pt
  }
}

% Box per il Prompt
\newtcolorbox{PromptBox}{
  chatstyle,
  colback=blue!5,
  colframe=blue!60!black,
  title=Prompt
}

% Box per ChatGPT corretto (verde)
\newtcolorbox{ChatGPTGood}{
  chatstyle,
  colback=green!5,
  colframe=green!60!black,
  title=Gemini
}

% Box per ChatGPT errato (rosso)
\newtcolorbox{ChatGPTBad}{
  chatstyle,
  colback=red!5,
  colframe=red!60!black,
  title=Gemini
}

\begin{document}

%TITLE PAGE

\begin{titlepage}

% --- BLOCCO SEGNALIBRO ---
\begin{tikzpicture}[remember picture, overlay]
    \coordinate (TopLeft) at (current page.north west);
    
    % --- MODIFICA QUI ---
    % Sostituisci '\draw' con '\fill'
    \fill[Blue, opacity=0.8] % <<< CAMBIATO QUI
        ([xshift=3cm]TopLeft) 
        -- ++(0, -6cm)        
        -- ++(1.25cm, 1.5cm)  
        -- ++(1.25cm, -1.5cm) 
        -- ++(0, 6cm)         
        -- cycle;             
\end{tikzpicture}
% --- FINE BLOCCO SEGNALIBRO ---


% --- INTESTAZIONE: INFO PROGETTO E LOGO ---
\noindent 
\begin{minipage}[t]{0.7\textwidth} 
    % Vuoto
\end{minipage}
\hfill 
\begin{minipage}[t]{0.45\textwidth} 
    \vspace{0pt} 
    \flushright 
    \includegraphics[width=\linewidth]{Images/PolimiLogo}
\end{minipage}

\vfill 

% --- TITOLO E SOTTOTITOLO ---
\begin{flushleft}
    {\textcolor{Blue}{\textbf{\Huge{Requirement Analysis and \\ Specification Document}}}}
    
    \vspace{1cm} 
    {\textcolor{MyGray}{\textbf{\LARGE{Best Bike Paths Project}}}}
\end{flushleft} 

\vfill

% --- PIÈ DI PAGINA: DATA ---
\begin{flushleft} 
    {\large 2025-2026}
\end{flushleft} 

\end{titlepage}
%Define deliverable specific info
%Replace cell contents where needed
\begin{table}[h!]
\begin{tabu} to \textwidth { X[0.3,r,p] X[0.7,l,p] }
\hline

\textbf{Deliverable:} & RASD\\
\textbf{Title:} & Requirement Analysis and Specification Document \\
\textbf{Authors:} & Andrea Ferritti, Jieke Lin, Tommaso Tosca \\
\textbf{Version:} & 1.0 \\ 
\textbf{Date:} & 15-November-2025 \\
\textbf{Download page:} & https://github.com/ferritti/FerrittiLinTosca \\
\textbf{Copyright:} & Copyright © 2025, A. Ferritti, J. Lin, T. Tosca – All rights reserved \\
\hline
\end{tabu}
\end{table}




\setcounter{page}{2}


%------------------------------------------------------------------------------------------------------------------------------------------------
\newpage
\addcontentsline{toc}{section}{Table of Contents}
\tableofcontents
\newpage
\addcontentsline{toc}{section}{List of Figures}
\listoffigures
\addcontentsline{toc}{section}{List of Tables}
\listoftables

%------------------------------------------------------------------------------------------------------------------------------------------------
\clearpage
{\color{Blue}{\section{Introduction}}}
\label{sect:introduction}
\subsection{Purpose}
The growing popularity of cycling as both a means of transportation and a recreational activity has increased the need for reliable, up-to-date information about bike paths. Cyclists often struggle to identify safe and efficient routes, especially when facing incomplete or outdated data regarding path conditions, obstacles, or environmental factors. At the same time, users who wish to contribute to the cycling community by sharing information about the paths they travel may find it difficult to organize, store, and communicate such data in a structured and meaningful way.
The purpose of BBP is to provide an environment where cyclists can record their trips, contribute information about bike paths, and access curated, community-driven data to make informed decisions about their routes. Through a combination of manual input and automated data collection from mobile devices, users can keep track of their cycling activities while simultaneously improving the shared knowledge base about path conditions.
BBP aims to give all users—whether registered or not—a clear understanding of the quality and safety of available bike paths, offering route visualizations, condition assessments, and recommendations based on consolidated and verified information. Transparency, reliability, and community collaboration are promoted at every step, ensuring that cyclists can navigate efficiently while benefiting from accurate and continuously updated path information.
\\

\noindent
BBP, therefore, intends to serve the following primary goals:

\begin{description}
    \item[\textbf{[G1]}] Registered users should have access to performance metrics about their trips.
    
    \item[\textbf{[G2]}] When available, users would like to be informed about meteorological information related to the paths of interest.
    
    \item[\textbf{[G3]}] Registered users should have the possibility to insert informations about the path of interest.
    
    \item[\textbf{[G4]}] Information about bike paths provided by registered users should be acquirable both manually and automatically.
    
    \item[\textbf{[G5]}] Automatically detected issues should be validated by the registered user before becoming available to others.
    
    \item[\textbf{[G6]}] Registered users should be able to make their provided information publishable to the community.
    
    \item[\textbf{[G7]}] Users should be able to visualize viable bike paths between two points.
    
    \item[\textbf{[G8]}] Users should be able to compare alternative bike paths based on their overall score.
\end{description}


\subsection{Scope}
The Best Bike Paths (BBP) system aims to support the biking community by mediating the interactions between individual cyclists and the shared repository of information regarding bike paths. The external users who interact with the platform include registered cyclists, unregistered visitors, and potentially local cycling associations or municipalities interested in monitoring the conditions of bike routes. Since only registered users can contribute new information or validate automatically collected data, the system must identify these users through personal accounts, while unregistered visitors access the platform in a read-only mode. Registered cyclists provide BBP with details about their trips either manually or through an automated acquisition mode that collects GPS data and movement information. In addition to recording the path itself, the application also provides users with performance metrics derived from their trips, such as total distance, duration, average speed, and other relevant indicators to support personal activity monitoring and analysis. These data sources become key assets for enhancing trip analysis, improving the accuracy of anomaly detection, and increasing the overall reliability of the platform’s scoring mechanisms.
When available, the system also enriches trip data with meteorological information, such as weather conditions, temperature, and wind speed, retrieved from external services.
The system supports users in reviewing, correcting, and publishing the collected information, and it aggregates contributions from multiple cyclists to ensure that the status of each path reflects fresh and trustworthy evidence. All external users can request the visualization of one or more paths between an origin and a destination, and BBP assists them by computing ranked paths based on their status and effectiveness. Through these features, the application facilitates the discovery of safe and efficient bike paths and provides a continuously updated overview of the cycling network.



\subsubsection{World Phenomena}
\begin{description}
    \item[\textbf{[WP1]}] A user moves along a path.
    
    \item[\textbf{[WP2]}] The user’s physical speed varies over time.
    
    \item[\textbf{[WP3]}] The user’s device experiences vibrations or impacts due to road conditions.
    
    \item[\textbf{[WP4]}] The user physically reaches the starting location of the trip.
    
    \item[\textbf{[WP5]}] The user arrives at the ending location of the trip.
    
    \item[\textbf{[WP6]}] The user trip is conditioned by the weather conditions.
\end{description}

\subsubsection{Shared Phenomena}
\paragraph{World Controlled}
\begin{description}
    \item[\textbf{[SP1]}]  The user publishes information about the path in manual mode

    \item[\textbf{[SP2]}] The user accepts path information suggested by automated mode

    \item[\textbf{[SP3]}] The system receives accelerometer readings generated by vibrations

    \item[\textbf{[SP4]}] The system receives gyroscope readings generated by device rotation 

    \item[\textbf{[SP5]}] The system receives GPS location data from the cyclist’s device

    \item[\textbf{[SP6]}]  Different users may observe the same path and report different aspects of it

    \item[\textbf{[SP7]}] The user inserts an origin and a destination

    \item[\textbf{[SP8]}] The user starts the navigation
\end{description}


\paragraph{Machine Controlled}
\begin{description}
    \item[\textbf{[SP9]}] The system suggests path information to be inserted in automated mode 

    \item[\textbf{[SP10]}] The system provides the top 5 paths based on the score

    \item[\textbf{[SP11]}] The system provides information about the path of interest

    \item[\textbf{[SP12]}] The system provides performance metrics about the user’s trip
\end{description}
\subsection{Definitions, Acronyms, Abbreviations}
\subsubsection{Definitions}

\begin{itemize}
    \item \textbf{Bike Path}: A road segment or sequence of segments that is suitable for bicycles. A path is considered bike-safe if it includes a dedicated bike lane or if car traffic is rare and speed limits are compatible with typical bike speeds.
    \item \textbf{Trip}: A cycling activity recorded by a registered user. A trip includes GPS tracking data, performance metrics (distance, average speed, etc.), and, when available, meteorological information retrieved from an external weather service.
    \item \textbf{Performance Metrics}: Statistics associated with a trip, such as total distance, duration, average speed, maximum speed, elevation gain, and other cycling-related measures.
    \item \textbf{Meteorological Information}: External data enriched by BBP when available, including temperature, weather conditions (sunny, cloudy, rainy…), wind speed, humidity, and similar environmental indicators.
    \item \textbf{Manual Mode}: A mode in which a user manually enters information about a bike path, specifying street names, path segments, their status, and any obstacles.
    \item \textbf{Automated Mode}: A mode in which BBP automatically collects data from the user’s mobile device while biking. Data includes GPS traces to infer the following path, accelerometer and gyroscope readings to detect potential potholes or obstacles, and speed information to infer that the user is cycling.
    \item \textbf{Obstacle}: Any relevant physical irregularity affecting the bike path (e.g., potholes, bumps, debris). In automated mode, obstacles are inferred from device sensors and later confirmed by the user.
    \item \textbf{Status}: A qualitative assessment of a bike path segment. Possible values include: optimal, medium, sufficient, requires maintenance, etc. Status can be user-generated (manual or automatic) and may be aggregated from multiple users.
    \item \textbf{Information Publishable}: Verified and user-approved data about a bike path (status, obstacles) that can be made publicly accessible to the community.
    \item \textbf{Score}: A computed score used to rank alternative bike paths between an origin and a destination. The score considers both the quality (status) of the path and its effectiveness in connecting the two points.
\end{itemize}

\subsubsection{Acronyms}
\begin{itemize}
    \item \textbf{BBP:} Best Bike Path.
    \item \textbf{GPS:} Global Positioning System.
\end{itemize}

\subsubsection{Abbreviations}
\begin{itemize}
    \item \textbf{G:} Goal.
    \item \textbf{SP:} Shared Phenomena.
    \item \textbf{WP: } World Phenomena.
\end{itemize}

\subsection{Revision history}
Version 1.0: 23/12/2025
\subsection{Reference Documents}
\begin{itemize}
    \item Assignment RDD AY 2025-2026
    \item Software Engineering 2 AY 2025/2026 slides "CreatingRASD"
\end{itemize}
\subsection{Document Structure}
\begin{enumerate}
    \item \textbf{Introduction:} It aims to provide a brief overview of the project, focusing in particular on the motivations behind it and the goals intended to be achieved through its development.
    \item \textbf{Overall Description:} It is a high-level description of how the system operates, including a detailed explanation of the phenomena involving the world, the machine, or both. This section also provides the domain description along with the related assumptions.
    \item \textbf{Specific Requirements:} This section presents a detailed analysis of the requirements needed to achieve the identified goals. In addition, it provides further information relevant to developers, such as hardware and software constraints.
    \item \textbf{Formal analysis:} It provides a formal specification of the world phenomena using Alloy.
    \item \textbf{Effort spent:} It reports the time spent to produce this document, organized by section.
    \item \textbf{References:}It includes references to all documents and software used in the preparation of this paper.
\end{enumerate}
 

%------------------------------------------------------------------------------------------------------------------------------------------------
\clearpage
{\color{Blue}{\section{Overall Description}}}
\label{sect:overview}
\subsection{Product perspective}

\subsubsection{Scenarios}

    \textbf{Registered users} 
    \vspace{2mm}
    \noindent


    \noindent
    \textbf{User publishes path information in automated mode} 
    \vspace{1mm}
    
    \noindent
    After completing the bike ride,  user U clicks the stop button related to the automated mode and will be redirected to the path data entry page with all the field 
    pre-filled by the system based on the data it collected, with the possibility to modify the status and the obstacles. Finally, U presses the "Publish" button, and the new path becomes immediately visible to all users.

    \vspace{3mm}
    \noindent
    \textbf{User publishes path information in manual mode} 
    \vspace{1mm}
    
    \noindent
    User U wants to share a path they frequently use with the community. From the main map interface, U presses the "Create Path" button (the plus icon). U enters the drawing mode and taps a sequence of points on the map to define the itinerary; the application automatically interprets these inputs and snaps the route line to the correct street geometry. If U is aware of specific obstacles along this route, such as deep potholes or dangerous intersections, U switches the tool to "Obstacles" mode and taps on the corresponding points of the path to flag them. Once the geometry is defined, U proceeds to the next step. The application presents a summary form where U enters the estimated duration and selects a status rating. The system automatically lists the names of the streets detected along the drawn path. Finally, U presses the "Publish" button, and the new path becomes immediately visible to all users.

    \vspace{3mm}
    \noindent
    \textbf{User eliminates a published path} 
    \vspace{1mm}
    
    \noindent
    User U wants to delete a previously published route. In the navbar, U clicks on “Published,” where he will find all the paths they have published. U clicks on the route they want to delete; the page for editing the route’s information will open, and they click on the icon associated with deleting the route.

    
    \vspace{3mm}
    \noindent
    \textbf{User edits the published path}
    \vspace{1mm}
    
    \noindent
    User U views their list of shared routes in the "Published" section and decides to update the details of a specific path due to changed road conditions. U taps on the path card to open the detailed view. To make changes, U presses the "Edit" button (the pencil icon) located in the header. The interface transitions to editing mode: the Start and End location names become editable text fields, and the read-only status badge is replaced by a selection grid, allowing U to quickly switch the rating (e.g., from "Maintenance Required" back to "Optimal"). Additionally, the description field becomes an active text area, allowing U to refine the route notes. Once the modifications are complete, U presses the "Save" button. The application updates the record and immediately reverts to the read-only detailed view, displaying the refreshed information.

    \vspace{3mm}
    \noindent
    \textbf{User modifies personal information}
    \vspace{1mm}
    
    \noindent
    User U wants to edit their personal information. In the navbar, U clicks on “Profile” and then on “Settings”, where he will find all their personal information with the option to modify it

    \newpage
    \noindent
    \textbf{User logs into the application}
    \vspace{1mm}
    
    \noindent
    User U, currently viewing the Home Page in the unregistered state, clicks on "My Profile" in the navbar. This action directs them to the Login screen. On this screen, U is prompted to enter their credentials: they fill in the email address in the designated field and the corresponding password in the password field. After ensuring the data is correct, U clicks on the "Login" button. Upon successful validation, the system grants U access and redirects them back to the main application Home Page, now displaying features and content available to a registered user.
    

    \vspace{3mm}
    \noindent
    \textbf{User discards the collected path information}
    \vspace{1mm}
    
    \noindent
    User U wants to discard the path information he entered. On the data entry screen for the given route, U clicks the “Discard” button, and the information will not be published to the community. However, if the information was entered in automated mode, it will still be added to the recent trips.

    \vspace{3mm}
    \noindent
    \textbf{User saves the trip}
    \vspace{1mm}
    
    \noindent
    User U wants to save the completed trip. Once the trip is finished, the end-of-path screen appears, showing the distance traveled and the average speed, along with the option to submit a report on the path’s status and obstacles. Finally, U clicks the "Back to Dashboard" button and the trips is saved into Recent Trips.


    \vspace{3mm}
    \noindent
    \textbf{User inserts an already existing path}
    \vspace{1mm}
    
    \noindent
    User u enters the information for a new path, but the streets in the entered route are the same as those of an existing path. In this case, the application will notify the user with a warning that the entered streets correspond to those of an already existing path. The same thing happens whether the insertion is done manually or automatically

    \vspace{3mm}

    \noindent
    \textbf{Unregistered users} 
    \vspace{2mm}
    
    \noindent
    \textbf{User registers in the application}
    \vspace{1mm}
    
    \noindent
    User U, having navigated to the Login screen and clicks the "Registration" link presented on this page, is immediately redirected to the dedicated registration page. Here, U is required to complete the registration form by filling in all the necessary information: they enter a valid email address and define a secure password. After reviewing the entered data, U presses the final "Register" button. A verification email will then be sent to the provided email address, containing a verification code that has to be inserted in the application in the appropriate textbox, to confirm that the entered email address exists. Upon successful completion of the registration process, the system confirms the creation of the new account and redirects U back to the main Login screen, prompting them to perform their first access.

    \vspace{3mm}

    \noindent
    \textbf{All users} 
    \vspace{2mm}
  
    \noindent
    \textbf{User searches for a path} 
    \vspace{1mm}
    
    \noindent
     User U wants to find the best path to their destination. From the main map interface, U enters the desired origin and destination points into the floating search bar. The system calculates and displays the five highest-scoring paths on the map. Each path is shown with its associated score, including the visualization of the corresponding obstacles present along the path. The path with the highest score is highlighted in teal to indicate the recommended route, while the other alternatives are displayed for comparison, allowing U to choose the best option before starting the ride.

     \vspace{2mm}
  
    \noindent
    \textbf{User starts a navigation} 
    \vspace{1mm}
    
    \noindent
     After entering the origin and destination and selecting the desired path, user U wants to start navigation.
    From the main map interface, U clicks the "Start Navigation" button, and the navigation screen opens. On the map, a cursor representing U’s current position is displayed; at the top, the current speed and street name are shown, while at the bottom there is an "Exit Navigation" button to end the navigation.

\newpage
\subsubsection{Domain Class Diagram}
    \begin{figure}[!h]
        \centering
        \includegraphics[width=1.1\linewidth]{RASD/Images/Domain Class Diagram.png}
        \label{fig:placeholder}
    \end{figure}

    \noindent
    Notes:
    \begin{itemize}
    \item The \textit{creates} relationship indicates the publication of the Path, while \textit{reports} tracks the contributions regarding obstacles and status, which is then aggregated by the \textit{System} (Merging).
    \item The \textit{status} attribute in \textit{Path} is initially set by the user who published the path, and is later derived from the aggregation of the most recent \textit{PathReport}.
    \item The \textit{structures} attribute in \textit{Path} is a list of points in the map that allows the \textit{System} to draw the path. In manual mode they are manually inserted on the map by the user who publishes the path, while in automated mode they are automatically inserted by the \textit{System}.
    \end{itemize}

\newpage
\subsubsection{State Diagram}
    \noindent
    \textbf{Trip lifecycle} 
    \vspace{2mm}
    \begin{figure}[!h]
        \centering
        \includegraphics[width=1\linewidth]{RASD/Images/SD_trip_lifecycle.png}
        \label{fig:placeholder}
    \end{figure}
\noindent \\
The Trip lifecycle state diagram illustrates the two operational modes through which the entity is managed by the system.In the upper branch, dedicated to assisted navigation, the user enters the desired origin and destination; the system identifies a set of paths with the highest scores, from which the user selects their preferred route. Navigation is then initiated: during this phase, the system guides the user to the destination while simultaneously recording driving performance metrics. The flow concludes with a review phase, during which data can still be modified before completion.In parallel, the lower branch describes the automated mode, designed for passive tracking: here, the trip immediately enters the recording phase, allowing for intermediate pauses. The recording must then be stopped to proceed to the final validation of the acquired data.Both flows finally converge in the concluding saving state, ensuring that, regardless of the input mode, all metrics, coordinates, and reports are correctly persisted in the system database.
    \newpage
    \vspace{2mm}
    \noindent
    \textbf{Path lifecycle} 
    \vspace{2mm}
    \begin{figure}[!h]
        \centering
        \includegraphics[width=1\linewidth]{RASD/Images/SD_path_lifecycle.png}
        \label{fig:placeholder}
    \end{figure}

\noindent
The Path lifecycle state diagram illustrates the process of creating and publishing a new path within the system. Starting from the initial state, the flow branches into two distinct operational modes.The manual mode allows the user to compose the path by interactively selecting existing streets of interest on the map and manually placing obstacle markers, with the ability to freely switch between these activities.Alternatively, the automated mode acquires the track via real-time GPS recording, handling temporary interruptions through the pause state.Both flows converge in the validation phase; in this state, the user can review and eventually modify the collected information before final confirmation. The cycle concludes with the publication of the path, or with the discard of the draft if not confirmed.

\newpage
    \vspace{2mm}
    \noindent
        \textbf{User Profile lifecycle} 
    \vspace{2mm}
    \begin{figure}[!h]
        \centering
        \includegraphics[width=1\linewidth]{RASD/Images/SD_user_profile_lifecycle.png}
        \label{fig:placeholder}
    \end{figure}

\noindent
The lifecycle begins when a visitor without an account completes registration, creating a persistent profile in the system that starts in a signed-out state. Entering credentials activates the working session, allowing the user to modify personal data or undertake new trips to improve their performance statistics. Upon finishing these activities, logging out returns the user to an inactive state while keeping the updated data saved for future access. The cycle ends permanently only if the user chooses to delete their account, an action that irreversibly removes the profile from the system.


\subsection{Product funtions}
    \noindent
    This section outlines the principal functionalities of the platform.

    \vspace{2mm}
    \noindent
    \textbf{Sign up \& Login} 
    \vspace{1mm}
    
    \noindent
     To sign up for the system, a user must provide his name, an email address, and a password. The system checks whether there is an existing account with the same email address: if so, the user is notified that the account already exists; otherwise, the registration proceeds and a verification code is sent to the provided email address. To complete the registration, the user must enter the verification code in the designated textbox.
     To log in, a user must enter the email and the associated password, and the system checks if the credentials match with an existing account: if so, the log in is completed; otherwise, the user is notified that the credentials are wrong.

    \vspace{2mm}
    \noindent
    \textbf{Path search and visualization} 
    \vspace{1mm}
    
    \noindent
     The system allows users to find the best paths that connect two locations. They must enter the starting point and destination, and the system will display on the map the paths with the highest score, along with their related obstacles. It is also possible to access the information screen of the selected path.

    \vspace{2mm}
    \noindent
    \textbf{Publication of a path in manual mode} 
    \vspace{1mm}
    
    \noindent
     The system allows registered users to publish new paths by manually providing the required information.
    The user draws the path and adds obstacles directly on the map, and then enters the remaining information.

    \vspace{2mm}
    \noindent
    \textbf{Publication of a path in automated mode} 
    \vspace{1mm}
    
    \noindent
     The system allows registered users to publish new paths by providing the required information automatically.
    The system retrieves data from GPS and sensors (accelerometer, gyroscope) while the user is biking, and automatically fills in all the necessary information, and the user can choose to accept, modify, or discard it.

    \vspace{2mm}
    \noindent
    \textbf{Navigation} 
    \vspace{1mm}
    
    \noindent
     After the user has selected the desired path from the map and starts navigation, the system provides a real-time representation of the user’s position on the map and also displays the current speed and the street being traveled.

    \vspace{2mm}
    \noindent
    \textbf{Trip tracking and analysis} 
    \vspace{1mm}
    
    \noindent
    The system allows registered users to keep track of their trips. After completing the trip, a screen will appear showing the trip statistics, and once the confirmation button is clicked, the system saves the trip and, if available, enriches it with meteorological information (weather, wind, temperature) retrieved from an external service.

    \vspace{2mm}
    \noindent
    \textbf{Path reporting} 
    \vspace{1mm}
    
    \noindent
    The system allows registered users to submit a report regarding the status and the presence of obstacles on an existing path.
    After completing a trip associated with the path in question, the user can modify the status and add or remove any obstacles, and once the confirmation button is clicked, the report will be sent to the system.

    \vspace{2mm}
    \noindent
    \textbf{Data Aggregation and Merging} 
    \vspace{1mm}
    
    \noindent
    The system collects path reports from multiple users and merges conflicting information. It calculates the current status of a path based on the freshness of the reports and the number of confirming data obtained.

\subsection{User characteristics}

    \subsubsection{Registered users}
    \noindent
    Registered users are able to login to the platform. Their profiles contain personal information and some statistics such as total kilometers travelled, total trips completed, total paths published and the average speed of their trips. In addition to all the features available to unregistered users, they can publish new paths (in either manual or automated mode) and submit reports regarding existing paths. Furthermore, they can record their trips to populate their history and view detailed statistics at the end of the navigation.

    \subsubsection{Unregistered users}
    \noindent
    Unregistered users include casual users, tourists, or people who simply want to find a bike path without tracking their data. They can search for the best paths between two locations, select a desired path, and start the real-time navigation to follow the path, but without the ability to save the trip or view post-activity statistics.

\newpage
\subsection{Assumptions, dependencies and constraints}

    \subsubsection{Domain Assumptions}
    \noindent
    This section describes the fundamental assumptions regarding the environment and the behavior of all entities interacting with the system. These assumptions simplify the system’s design and implementation by setting clear expectations for the users, and the platform itself, ensuring the system functions efficiently within its intended contex

    \begin{description}
        \item[\textbf{[D1]}] The user’s device is equipped with a GPS, an accelerometer and a gyroscope.
        \item[\textbf{[D2]}] The user has a working internet connection.
        \item[\textbf{[D3]}] The user provides truthful information regarding the path to be published.
        \item[\textbf{[D4]}] The user provides truthful reports regarding paths.
        \item[\textbf{[D5]}] The external weather information service is reliable.
        \item[\textbf{[D6]}] The external map service is available and provides accurate topographical data.
    \end{description}

    \subsubsection{Dependencies}
    \noindent
    \begin{itemize}
    \item \textbf{Map Service}: it's responsible for providing the map that enables the graphical visualization of the paths. This is a critical dependency: if this service is unavailable, the core functionalities of the system cannot be performed.
    \item \textbf{Weather Service}: it's responsible for adding weather information to the users’ trips. This is a non-critical dependency: if this service is unavailable, the system will continue to function normally, simply omitting the weather information from the trip details. 
    \end{itemize}

%------------------------------------------------------------------------------------------------------------------------------------------------
\clearpage
{\color{Blue}{\section{Specific Requirements}}}
\label{sect:requirements}
\subsection{External Interface Requirements}
Here is shown the mockups of the UI. With the exception of the home page, which is shown from both the guest’s and the user’s perspectives, all other screens are displayed only from the user’s perspective, since the only differences are already shown on the home page.
\vspace{3mm}
\newline
\textbf{Home page} \newline
The \emph{Home page} acts as a dashboard for exploration. It assumes the cyclist immediately wants to either find a safe path to ride or start recording their own journey. On the left, the home page is shown from the user’s perspective; on the right, from the guest’s perspective

\begin{figure}[H]
    \centering
    \includegraphics[width=1\linewidth]{RASD/Images/Mockup/HomePage.png}
    \caption{Home page}
    \label{fig:placeholder}
\end{figure}

\newpage
\noindent
\textbf{Search path} \newline
From the \emph{Home page}, when the cyclist inserts a destination, the five paths that connect origin and destination with the highest score will appear on the map. If the cyclist wants to change the origin to a location that is not the current location, the “Start Navigation” button will disappear.

\begin{figure}[H]
    \centering
    \includegraphics[width=0.6\linewidth]{RASD/Images/Mockup/StartNavigation.png}
    \caption{Start navigation screen}
    \label{fig:placeholder}
\end{figure}

\newpage
\noindent
\textbf{Navigation} \newline
From the \emph{Search path screen}, when the cyclist clicks “Start Navigation”, the navigation begin and the \emph{Navigation sreen} appears. When the cyclist reaches the destination, the navigation ends and the \emph{Trip completed screen} appears.

\begin{figure}[h!]
    \centering
    \subfloat[Navigation screen]{\includegraphics[width=0.5\textwidth]{RASD/Images/Mockup/Navigation.png}}\hfill
    \subfloat[Trip completed screen]{\includegraphics[width=0.5\textwidth]{RASD/Images/Mockup/NavigationCompleted.png}}
\end{figure}

\newpage
\noindent
\textbf{Manual Mode} \newline
From the \emph{Home page}, when the user clicks the “Manual Mode” button, the \emph{Draw Path screen} appears. Here, the user can draw a path by clicking on the map to add points, which the Map Service then connects to create a path that follows the street layout. By clicking “Obstacles”, the \emph{Insert Obstacles screen} appears, allowing the user to place obstacles directly on the map. Finally, clicking “Next Step” opens the \emph{Review \& Publish screen}, where the user sets the status and publishes the path by clicking “Publish Path”.
\begin{figure}[h!]
    \centering
    \subfloat[Draw path screen]{\includegraphics[width=0.33\textwidth]{RASD/Images/Mockup/DrawPath.png}}\hfill
    \subfloat[Insert obstacles screen]{\includegraphics[width=0.33\textwidth]{RASD/Images/Mockup/InsertObstacles.png}}
    \subfloat[Review \& publish screen]{\includegraphics[width=0.33\textwidth]{RASD/Images/Mockup/Review&Publish.png}}
\end{figure}

\newpage
\noindent

\newpage
\subsection{Functional Requirements}

\textbf{Sign up \& Login}
\begin{description}
    \item[\textbf{[R1]}] The platform shall allow new users to register an account by providing personal information.
    \item[\textbf{[R2]}] The platform shall allow registered users to log in by entering their registered email and password.
\end{description}

\noindent
\textbf{Map Explorer}
\begin{description}
    \item[\textbf{[R3]}] The platform shall display an interactive map centered on the user's current location.
    \item[\textbf{[R4]}] The platform shall allow users to input an "Origin" and "Destination" of the desired path.
    \item[\textbf{[R5]}] The platform shall allow users to swap the origin and destination inputs.
    \item[\textbf{[R6]}] The platform shall calculate and display the top 5 paths options ranked by score.
    \item[\textbf{[R7]}] The platform shall visualize the geometry, quality, scores and obstacles of the top 5 paths on the map.
    \item[\textbf{[R8]}] The platform shall allow users to start the navigation mode for the currently selected path.
    \item[\textbf{[R9]}] When the navigation is started, the platform shall visualize the user's position in real-time.
    \item[\textbf{[R10]}] When the navigation is started, the platform shall display a dynamic navigation overlay providing the name of the street being traversed and the estimated speed.

    \item[\textbf{[R11]}] the platform shall allow users to interrupt the active navigation mode at any time.

    \item[\textbf{[R12]}]  When the navigation is completed, the platform shall present to registered users a summary screen containing statistics and the map of the path taken.

    \item[\textbf{[R13]}] When the navigation is completed, the platform shall allow registered users to modify the path status and the obstacles.
\end{description}

\noindent
\textbf{Active Path Recording}
\begin{description}
    \item[\textbf{[R14]}] The platform shall allow registered users to start an active recording session.
    \item[\textbf{[R15]}] The platform shall allow registered users to pause and resume the recording of the active trip.
    \item[\textbf{[R16]}] The platform shall utilize device sensors to automatically detect obstacles.
    \item[\textbf{[R17]}] The platform shall display real-time metrics during the trip, including duration, current speed, distance traveled, and weather conditions.

    \item[\textbf{[R18]}] The platform shall allow registered users to stop the recording and immediately transition to the summary screen to review the captured data.

    \item[\textbf{[R19]}] The platform shall allow registered users to modify status and obstacles in the captured data summary.

    \item[\textbf{[R20]}] The platform shall allow registered users to publish the captured path along with its data.
\end{description}

\newpage
\noindent
\textbf{Manual Path Creation}
\begin{description}
    \item[\textbf{[R21]}] The platform shall allow registered users to manually draw the path on the map.
    \item[\textbf{[R22]}] The platform shall allow registered users to manually place obstacle markers along the drawn path.
    \item[\textbf{[R23]}] The platform shall allow registered users to set the status of the drawn path.
    \item[\textbf{[R24]}] The platform shall allow registered users to publish the drawn path along with its data.
\end{description}

\noindent
\textbf{Trips \& Published Paths}
\begin{description}
    \item[\textbf{[R25]}] The platform shall display a dashboard listing the user's historical trips and their data, sorted chronologically.
    \item[\textbf{[R26]}] The platform shall allow registered users to view a list of paths published by them.
    \item[\textbf{[R27]}] The platform shall allow registered users to edit the information of any of their published paths.
    \item[\textbf{[R28]}] The platform shall allow registered users to delete any of their published paths.
\end{description}

\noindent
\textbf{User Profile}
\begin{description}
    \item[\textbf{[R29]}] The platform shall display a user profile screen summarizing the main aggregate statistics.
    \item[\textbf{[R30]}] The platform shall allow registered users to edit their profile information.
\end{description}



\newpage
\subsubsection{Use Cases Diagrams}
    \begin{figure}[H]
        \centering
        \includegraphics[width=1\linewidth]{RASD/Images/UseCaseDiagrams/GuestCyclistUser.jpg}
        \caption{General Utilities Use Case Diagram}
        \label{fig:placeholder}
    \end{figure}

This Use Case Diagram depicts the system’s functionalities organized through a generalization hierarchy. The abstract actor Cyclist represents the baseline capability of searching for paths, which is common to all users. Two actors inherit from this base:

\begin{itemize}
    \item The Guest (unauthenticated), who extends the base capability with access-control operations.
    \item The User (authenticated), who extends the base capability with the full suite of platform features.
\end{itemize}

Finally, the diagram illustrates the participation of external services (Map and Weather Services) in supporting these functionalities.
\newpage
\begin{figure}[H]
    \centering
    \includegraphics[width=1\linewidth]{RASD/Images/UseCaseDiagrams/ManualMode.jpg}
    \caption{Manual Mode Use Case Diagram}
    \label{fig:manual_mode}
\end{figure}

\noindent
This diagram illustrates the manual path creation process centered on the 'Publish Path' use case. This main activity mandatorily includes both 'Draw Path' (to define the route geometry) and 'Add Status' (to assign a rating). Additionally, the process can be optionally extended by 'Add Obstacle' to allow the user to report specific issues on the map.

\vspace{15mm}

\begin{figure}[H]
    \centering
    \includegraphics[width=1\linewidth]{RASD/Images/UseCaseDiagrams/AutomatedMode.jpg}
    \caption{Automated Mode Use Case Diagram}
    \label{fig:placeholder}
\end{figure}
\noindent
This diagram illustrates the automated path acquisition process. The Device Sensors actor refers to the device’s GPS, accelerometer, and gyroscope, which are essential for tracking the user's movement and detecting path quality in real-time.
\newpage

\subsubsection{Use Cases}

\setcounter{table}{0}
%tabella
\begin{table}[H]
    \centering
    \setlength\extrarowheight{2pt}
    \renewcommand{\arraystretch}{1.2}
    \begin{tabularx}{\textwidth}{|>{\bfseries}l|X|}
    \hline
    UC-01 & Sign up \\ \hline
    Actors & Guest\\ \hline
    Entry Condition &  The Guest is not currently logged in.\\ \hline
    Event Flow & 
    \begin{enumerate}
      \item The Guest opens the login screen.
      \item The Guest clicks "sign up".
      \item The Guest enters the required information.
      \item The Guest presses the "sign up" button.
      \item The platform validates the entered information.
      \item The platform sends a confirmation email to the provided email address.
      \item The platform shows a message to the guest to confirm the registration.
      \item The Guest confirms the registration by clicking on the link in the email. 
      \item The platform confirms the registration and activates the account.
      \item The Guest is redirected to the platform's homepage as a logged-in User.
    \end{enumerate}
    \\ \hline
    Exit Condition &  A new account is created and the Guest is now logged in as a User.\\ \hline
    Exception & 
    \begin{itemize}
        \item The Guest provides incorrect information.
        \item The Guest does not confirm the registration.
        \item The email is already in use.
    \end{itemize}
    \\ \hline
    \end{tabularx}
    \caption{Use Case UC-01}
    \label{UC-01}
\end{table}

\newpage
\begin{table}[H]
    \centering
    \setlength\extrarowheight{2pt}
    \renewcommand{\arraystretch}{1.2}
    \begin{tabularx}{\textwidth}{|>{\bfseries}l|X|}
    \hline
    UC-02 & Login \\ \hline
    Actors & Guest\\ \hline
    Entry Condition & The Guest is not currently logged in.\\ \hline
    Event Flow & 
    \begin{enumerate}
      \item The Guest opens the login screen.
      \item The Guest enters login credentials.
      \item The Guest clicks the "login" button.
      \item The platform verifies login credentials.
      \item The platform creates a login session for the guest.
      \item The Guest is redirected to the platform’s homepage as a logged-in User.
    \end{enumerate}
    \\ \hline
    Exit Condition & The Guest is successfully authenticated as a User.\\ \hline
    Exception & 
    \begin{itemize}
        \item The credentials provided are incorrect.
    \end{itemize}
    \\ \hline
   \end{tabularx}
   \caption{Use Case UC-02}
    \label{UC-02}
\end{table}

\begin{table}[H]
    \centering
    \setlength\extrarowheight{2pt}
    \renewcommand{\arraystretch}{1.2}
    \begin{tabularx}{\textwidth}{|>{\bfseries}l|X|}
    \hline
    UC-03 & Search Path \\ \hline
    Actors & Cyclist, Map Service\\ \hline
    Entry Condition & The Cyclist is viewing the platform’s homepage\\ \hline
    Event Flow & 
    \begin{enumerate}
      \item The platform loads the visual map from the Map Service.
      \item The Cyclist enters the origin.
      \item The Cyclist enters the destination.
      \item The platform retrieves the locations using the Map Service.
      \item The platform identifies the best 5 paths among those published.
      \item The platform displays the best 5 paths on the map.
      \item The Cyclist selects the desired path.
    \end{enumerate}
    \\ \hline
    Exit Condition &  The desired path is selected and the option to view its details is available to the Cyclist.\\ \hline
    Exception & 
    \begin{itemize}
        \item None.
    \end{itemize}
    \\ \hline
    \end{tabularx}
    \caption{Use Case UC-03}
    \label{UC-03}
\end{table} 

\begin{table}[H]
    \centering
    \setlength\extrarowheight{2pt}
    \renewcommand{\arraystretch}{1.2}
    \begin{tabularx}{\textwidth}{|>{\bfseries}l|X|}
    \hline
    UC-04 & Start Navigation \\ \hline
    Actors & Cyclist\\ \hline
    Entry Condition & The Cyclist has selected the desired path and the origin is the current location\\ \hline
    Event Flow & 
    \begin{enumerate}
      \item The platform shows the "Start Navigation" button.
      \item The Cyclist clicks "Start Navigation".
      \item The platform displays the selected path on the map.
      \item The platform guides the Cyclist in real-time.
      \item The platform detects that the Cyclist has reached the destination.
    \end{enumerate}
    \\ \hline
    Exit Condition &  The platform displays the trip summary.\\ \hline
    Exception & 
    \begin{itemize}
        \item None.
    \end{itemize}
    \\ \hline
   \end{tabularx}
   \caption{Use Case UC-04}
    \label{UC-04}
\end{table}

\begin{table}[H]
    \centering
    \setlength\extrarowheight{2pt}
    \renewcommand{\arraystretch}{1.2}
    \begin{tabularx}{\textwidth}{|>{\bfseries}l|X|}
    \hline
    UC-05 & Save trip \\ \hline
    Actors & User, Weather Service\\ \hline
    Entry Condition & The user is logged in and the navigation is completed.\\ \hline
    Event Flow & 
    \begin{enumerate}
      \item The platform retrieves weather data from the Weather Service.
      \item The platform displays the "Trip Completed" screen with statistics, weather info, and rating options. 
      \item The User selects the path status.
      \item The User clicks "Modify Obstacles" to report obstacles on the map.
      \item The User clicks "Save Trip".
      \item The platform saves the trip in the User's "Recent Trips" list.
      \item The platform automatically submits the path report to the system.
      \item The User is redirected to the "Recent Trips" screen.
    \end{enumerate}
    \\ \hline
    Exit Condition & The trip is saved and the "Recent Trips" screen is displayed to the User.\\ \hline
    Exception & 
    \begin{itemize}
        \item None.
    \end{itemize}
    \\ \hline
   \end{tabularx}
   \caption{Use Case UC-05}
    \label{UC-05}
\end{table}

\begin{table}[H]
    \centering
    \setlength\extrarowheight{2pt}
    \renewcommand{\arraystretch}{1.2}
    \begin{tabularx}{\textwidth}{|>{\bfseries}l|X|}
    \hline
    UC-06 & View Trips \\ \hline
    Actors & User\\ \hline
    Entry Condition & The User is logged in.\\ \hline
    Event Flow & 
    \begin{enumerate}
      \item The User opens the "Recent Trips" screen.
      \item The platform displays the list of saved trips.
      \item The User views the details of the trips in the list.
    \end{enumerate}
    \\ \hline
    Exit Condition &  The detailed view of the saved trips is displayed.\\ \hline
    Exception & 
    \begin{itemize}
        \item None.
    \end{itemize}
    \\ \hline
   \end{tabularx}
   \caption{Use Case UC-06}
    \label{UC-06}
\end{table}

\begin{table}[H]
    \centering
    \setlength\extrarowheight{2pt}
    \renewcommand{\arraystretch}{1.2}
    \begin{tabularx}{\textwidth}{|>{\bfseries}l|X|}
    \hline
    UC-07 & Edit Profile \\ \hline
    Actors & User\\ \hline
    Entry Condition & The User is logged in.\\ \hline
    Event Flow & 
    \begin{enumerate}
      \item The User opens the "Profile" screen.
      \item The User clicks on "Edit Profile".
      \item The User modifies the information.
      \item The User clicks "Confirm Changes"
      \item The platform validates the new data.
      \item The platform updates the profile and displays the success message.
    \end{enumerate}
    \\ \hline
    Exit Condition &  The User profile is updated with the new information.\\ \hline
    Exception & 
    \begin{itemize}
        \item The User leaves a mandatory field empty and clicks "Confirm"
    \end{itemize}
    \\ \hline
   \end{tabularx}
   \caption{Use Case UC-07}
    \label{UC-07}
\end{table}

\begin{table}[H]
    \centering
    \setlength\extrarowheight{2pt}
    \renewcommand{\arraystretch}{1.2}
    \begin{tabularx}{\textwidth}{|>{\bfseries}l|X|}
    \hline
    UC-08 & Delete path \\ \hline
    Actors & User\\ \hline
    Entry Condition & The User is logged in.\\ \hline
    Event Flow & 
    \begin{enumerate}
      \item The User opens the "Published path" screen.
      \item The User selects the path to delete.
      \item The platform displays the path details.
      \item The User clicks the "Delete" button.
      \item The platform asks for confirmation.
      \item The User confirms the deletion.
      \item The platform removes the path and updates the list.
    \end{enumerate}
    \\ \hline
    Exit Condition &  The path is permanently deleted from the published paths list.\\ \hline
    Exception & 
    \begin{itemize}
        \item None.
    \end{itemize}
    \\ \hline
   \end{tabularx}
   \caption{Use Case UC-08}
    \label{UC-08}
\end{table}

\begin{table}[H]
    \centering
    \setlength\extrarowheight{2pt}
    \renewcommand{\arraystretch}{1.2}
    \begin{tabularx}{\textwidth}{|>{\bfseries}l|X|}
    \hline
    UC-09 & Publish Path (Manual Mode) \\ \hline
    Actors & User\\ \hline
    Entry Condition & The User is logged in and is viewing the homepage.\\ \hline
    Event Flow & 
    \begin{enumerate}
      \item The User clicks the "+" button to create a new path.
      \item The platform loads the "Draw Mode" interface.
      \item The User draw the path.
      \item The User clicks the "Next Step" button.
      \item The platform displays the "Review \& Publish" screen.
      \item The User selects the status.
      \item The User clicks "Publish Path".
      \item The platform validates and makes the path available to the community.
    \end{enumerate}
    \\ \hline
    Exit Condition &  The path is published and appears in the User's "Published Paths" screen.\\ \hline
    Exception &
    \begin{itemize}
        \item  The platform detects that an identical path already exists.
    \end{itemize}
    \\ \hline
   \end{tabularx}
   \caption{Use Case UC-09}
    \label{UC-09}
\end{table}

\begin{table}[H]
    \centering
    \setlength\extrarowheight{2pt}
    \renewcommand{\arraystretch}{1.2}
    \begin{tabularx}{\textwidth}{|>{\bfseries}l|X|}
    \hline
    UC-10 & Draw Path \\ \hline
    Actors & User, Map Service\\ \hline
    Entry Condition & The User is in the 'Draw Path' mode within the manual creation interface and the interactive map is loaded.\\ \hline
    Event Flow & 
    \begin{enumerate}
      \item The User taps on the map to place the starting point.
      \item The User continues tapping to trace the path segments.
      \item The Map Service renders the route line in real-time.
      \item The platform calculates the "Total Km" dynamically.
      \item The User completes the tracing.
    \end{enumerate}
    \\ \hline
    Exit Condition &  The geometric path is defined on the map.\\ \hline
    Exception & 
    \begin{itemize}
        \item The platform fails to load the map tiles from the Map Service.
    \end{itemize}
    \\ \hline
   \end{tabularx}
   \caption{Use Case UC-10}
    \label{UC-10}
\end{table}

\begin{table}[H]
    \centering
    \setlength\extrarowheight{2pt}
    \renewcommand{\arraystretch}{1.2}
    \begin{tabularx}{\textwidth}{|>{\bfseries}l|X|}
    \hline
    UC-11 & Add Obstacle \\ \hline
    Actors & User, Map Service\\ \hline
    Entry Condition & The User is in the manual creation interface, the interactive map is loaded, and the path is partially or fully drawn.\\ \hline
    Event Flow & 
    \begin{enumerate}
      \item The User clicks on the "Obstacles" tab.
      \item The platform switches to "Obstacle Mode".
      \item The User taps a specific point along the drawn path.
      \item The platform places an obstacle marker on the coordinate using the Map Service.
    \end{enumerate}
    \\ \hline
    Exit Condition & An obstacle is flagged on the map.\\ \hline
    Exception & 
    \begin{itemize}
        \item The user adds an obstacle outside the path
    \end{itemize}
    \\ \hline
   \end{tabularx}
   \caption{Use Case UC-11}
    \label{UC-11}
\end{table}

\begin{table}[H]
    \centering
    \setlength\extrarowheight{2pt}
    \renewcommand{\arraystretch}{1.2}
    \begin{tabularx}{\textwidth}{|>{\bfseries}l|X|}
    \hline
    UC-12& Start Automated Path Acquisition \\ \hline
    Actors & User, Device Sensors\\ \hline
    Entry Condition & The User is logged in and is viewing the homepage\\ \hline
    Event Flow & 
    \begin{enumerate}
      \item The User clicks the "Lightning" button to start the acquisition.
      \item The platform initializes the "Automated Mode" interface.
      \item The platform loads the User's current location on the map using the gps from the Device Sensors.
      \item The platform starts the timer.
      \item The platform begins recording distance, speed and obstacles in real-time using the Device Sensors.
      \item The platform continuously updates the metrics (distance, obstacles, speed).
    \end{enumerate}
    \\ \hline
    Exit Condition & The path recording is active and the "Automated Mode" interface is displayed.\\ \hline
    Exception & 
    \begin{itemize}
        \item None
    \end{itemize}
    \\ \hline
   \end{tabularx}
   \caption{Use Case UC-12}
    \label{UC-12}
\end{table}

\begin{table}[H]
    \centering
    \setlength\extrarowheight{2pt}
    \renewcommand{\arraystretch}{1.2}
    \begin{tabularx}{\textwidth}{|>{\bfseries}l|X|}
    \hline
    UC-13& Pause Path Acquisition \\ \hline
    Actors & User\\ \hline
    Entry Condition & The User is logged in and the "Automated Mode" is active\\ \hline
    Event Flow & 
    \begin{enumerate}
      \item The User clicks the "Pause" button.
      \item The platform suspends the timer.
      \item The platform stops recording coordinates and metrics.
      \item The platform processes the information.
      \item The platform updates the interface to show the "Paused" state.
    \end{enumerate}
    \\ \hline
    Exit Condition & The "Automated Mode" is temporarily suspended.\\ \hline
    Exception & 
    \begin{itemize}
        \item None
    \end{itemize}
    \\ \hline
   \end{tabularx}
   \caption{Use Case UC-13}
    \label{UC-13}
\end{table}

\begin{table}[H]
    \centering
    \setlength\extrarowheight{2pt}
    \renewcommand{\arraystretch}{1.2}
    \begin{tabularx}{\textwidth}{|>{\bfseries}l|X|}
    \hline
    UC-14& Stop Path Acquisition \\ \hline
    Actors & User, Map Service\\ \hline
    Entry Condition & The User is logged in and the "Automated Mode" is active (or paused)\\ \hline
    Event Flow & 
    \begin{enumerate}
      \item The User clicks the "Stop" button.
      \item The platform stops the timer.
      \item The platform stops the data recording.
      \item The platform processes the recorded data to calculate final metrics.
      \item The platform extracts the street names using the Map Service.
      \item The platform displays the "Trip Summary" interface.
    \end{enumerate}
    \\ \hline
    Exit Condition &  The "Automated Mode" is stopped and the User is viewing the "Trip Summary"\\ \hline
    Exception & 
    \begin{itemize}
        \item None
    \end{itemize}
    \\ \hline
   \end{tabularx}
   \caption{Use Case UC-14}
    \label{UC-14}
\end{table}

\begin{table}[H]
    \centering
    \setlength\extrarowheight{2pt}
    \renewcommand{\arraystretch}{1.2}
    \begin{tabularx}{\textwidth}{|>{\bfseries}l|X|}
    \hline
    UC-15& Publish Path (Automated Mode) \\ \hline
    Actors & User\\ \hline
    Entry Condition & The User is viewing the "Trip Summary" interface with the recorded path data.\\ \hline
    Event Flow & 
    \begin{enumerate}
      \item The User clicks the "Publish Path" button.
      \item The Platform stores the path data (geometry, status, obstacles) in the system.
      \item The platform validates and makes the path available to the community..
    \end{enumerate}
    \\ \hline
    Exit Condition &  The path is published and appears in the User’s "Published Paths" screen.\\ \hline
    Exception & 
    \begin{itemize}
        \item The platform detects that an identical path already exists.
    \end{itemize}
    \\ \hline
   \end{tabularx}
   \caption{Use Case UC-15}
    \label{UC-15}
\end{table}


\newpage
\subsubsection{Sequence Diagrams}
The functions used in the sequence diagrams follow these conventional meanings:
\begin{itemize}
    \item \textbf{press(Button)}: Represents the action of a user pressing Button.
    \item \textbf{submit(Content)}: Indicates a user uploading Content to the platform.
    \item \textbf{showPage(Page)}: Displays the content of a specific Page.
    \item \textbf{showError(Error)}: Displays an error message indicating an issue that occurred.
    \item \textbf{check()}: Verifies the validity of the information provided by the user.
    \item \textbf{update(Data)}: Refers to internal processes that update database records with Data.
    \item \textbf{OpenApplication()}: Rappresent the action of a user opens the application.
    \item \textbf{getMapData()}: Represents the system requesting the visual map components from the external service.
    \item \textbf{mapTiles(Tiles)}: Delivers the graphical data required to render the map interface on the screen.
    \item \textbf{enter(Location)}: Represents the action of a user inputting a starting point or destination address.
    \item \textbf{findBestPaths(Origin, Destination)}: Identifies and returns the 5 paths with the best performance scores connecting the selected locations.
    \item \textbf{selectPath(Path)}: Represents the user choosing one of the displayed paths to view its details or start navigation.
    \item \textbf{startTracking()}: Commands the device sensors (GPS) to begin sending real-time location updates.
    \item \textbf{stopTracking()}: Commands the device sensors to cease location updates
    \item \textbf{getWeather(CurrentLocation)}: Requests meteorological data (temperature, conditions) for the specific coordinates where the trip ended.
    \item \textbf{wetherData()}: Returns the requested weather information from the external service.
    \item \textbf{insert(Object)}: Represents the user placing a marker on the map.
    \item \textbf{initializeSensors()}: Internal process that activates and configures the device's hardware to start receiving real-time data.
    \item \textbf{monitorRoadConditions()}: Analyzes data from the accelerometer and gyroscope to automatically detect potholes during the ride.
    \item \textbf{displayInstruction()}: Updates the user interface with visual navigation cues and the current position on the map.
    \item \textbf{fetchData(Criteria)}: A generic system request to retrieve specific datasets from the database required to populate or plot elements on the user interface.
    \item \textbf{showMessage(Argument)}: Shows a general message on the screen.
    \item \textbf{startTracking()}: Activates the continuous recording and initiates the loop.
    \item \textbf{monitorRoadConditions()}:Processes live sensor data to understand the path environment.
    \item \textbf{stopTracking()}: Commands the device sensors to cease location updates.

    
\end{itemize}
\newpage
\vspace{2mm}
\noindent
    \textbf{[SD1] Sign up Sequence Diagram} 
\vspace{2mm}
\begin{figure}[!h]
    \centering
    \includegraphics[width=0.8\linewidth]{RASD/Images/SEQD_SignUp.png}
    \label{fig:placeholder}
\end{figure}
\newpage
\vspace{2mm}
\noindent
    \textbf{[SD2] Login Sequence Diagram} 
\vspace{2mm}
\begin{figure}[!h]
    \centering
    \includegraphics[width=0.8\linewidth]{RASD/Images/SEQD_Login.png}
    \label{fig:placeholder}
\end{figure}
\newpage
\vspace{2mm}
\noindent
    \textbf{[SD3] Search Path Sequence Diagram} 
\vspace{2mm}
\begin{figure}[!h]
    \centering
    \includegraphics[width=0.8\linewidth]{RASD/Images/SEQD_search_path.png}
    \label{fig:placeholder}
\end{figure}
\newpage
\vspace{2mm}
\noindent
    \textbf{[SD4] Start Navigation Sequence Diagram} 
\vspace{2mm}
\begin{figure}[!h]
    \centering
    \includegraphics[width=0.8\linewidth]{RASD/Images/SEQD_start_navigation.png}
    \label{fig:placeholder}
\end{figure}
\vspace{2mm}
\newpage
\noindent
    \textbf{[SD5] Save Trip Sequence Diagram} 
\vspace{2mm}
\begin{figure}[!h]
    \centering
    \includegraphics[width=0.8\linewidth]{RASD/Images/SEQD_save_trip.png}
    \label{fig:placeholder}
\end{figure}
\newpage
\noindent
    \textbf{[SD6] Trip View Sequence Diagram} 
\vspace{2mm}
\begin{figure}[!h]
    \centering
    \includegraphics[width=0.8\linewidth]{RASD/Images/SEQD_view_trip.png}
    \label{fig:placeholder}
\end{figure}
\newpage
\noindent
    \textbf{[SD7] Edit Profile Sequence Diagram} 
\vspace{2mm}
\begin{figure}[!h]
    \centering
    \includegraphics[width=0.8\linewidth]{RASD/Images/SEQD_edit_profile.png}
    \label{fig:placeholder}
\end{figure}
\newpage
\noindent
    \textbf{[SD8] Delete Path Sequence Diagram} 
\vspace{2mm}
\begin{figure}[!h]
    \centering
    \includegraphics[width=0.8\linewidth]{RASD/Images/SEQD_delete_path.png}
    \label{fig:placeholder}
\end{figure}
\newpage
\noindent
    \textbf{[SD9] Publish Path Sequence Diagram} 
\vspace{2mm}
\begin{figure}[!h]
    \centering
    \includegraphics[width=0.8\linewidth]{RASD/Images/SEQD_publish_path_manual_mode.png}
    \label{fig:placeholder}
\end{figure}
\newpage
\noindent
    \textbf{[SD10] Draw Path Sequence Diagram} 
\vspace{2mm}
\begin{figure}[!h]
    \centering
    \includegraphics[width=0.8\linewidth]{RASD/Images/SEQD_draw_path.png}
    \label{fig:placeholder}
\end{figure}
\newpage
\noindent
    \textbf{[SD11] Add Obstacle Sequence Diagram} 
\vspace{2mm}
\begin{figure}[!h]
    \centering
    \includegraphics[width=0.8\linewidth]{RASD/Images/SEQD_add_obstacle.png}
    \label{fig:placeholder}
\end{figure}
\newpage
\noindent
    \textbf{[SD12] Start Automated Path Acquisition Sequence Diagram} 
\vspace{2mm}
\begin{figure}[!h]
    \centering
    \includegraphics[width=0.8\linewidth]{RASD/Images/SEQD_start_automated_path_acquisition.png}
    \label{fig:placeholder}
\end{figure}
\newpage
\noindent
    \textbf{[SD13] Pause Path Acquisition Sequence Diagram} 
\vspace{2mm}
\begin{figure}[!h]
    \centering
    \includegraphics[width=0.8\linewidth]{RASD/Images/SEQD_pause_path_acquisition.png}
    \label{fig:placeholder}
\end{figure}
\newpage
\noindent
    \textbf{[SD14] Stop Path Acquisition Sequence Diagram} 
\vspace{2mm}
\begin{figure}[!h]
    \centering
    \includegraphics[width=0.8\linewidth]{RASD/Images/SEQD_stop_path_acquisition.png}
    \label{fig:placeholder}
\end{figure}
\newpage
\noindent
    \textbf{[SD15] Publish Path Automatated Mode Sequence Diagram} 
\vspace{2mm}
\begin{figure}[!h]
    \centering
    \includegraphics[width=0.8\linewidth]{RASD/Images/SEQD_publish_path_automated_mode.png}
    \label{fig:placeholder}
\end{figure}

\newpage
\subsection{Requirements Mapping}
\begin{table}[H]
    \centering
    % Aumenta lo spazio tra le righe per leggibilità
    \renewcommand{\arraystretch}{1.5}
    
    % Definizione tabella: larghezza totale testo, bordo verticale, colonna stretta per ID, colonna larga X per testo
    \begin{tabularx}{\textwidth}{|l X|}
        \hline
        % Riga Intestazione (G1) che occupa due colonne
        \multicolumn{2}{|l|}{\textbf{G1}: Registered users should have access to performance metrics about their trips.} \\
        \hline
        & \\[-1em] % Piccolo aggiustamento spaziatura
        
        % Domain Assumptions
        {[D12]} & The user’s device is equipped with a GPS, an accelerometer and a gyroscope. \\
        {[D2]} & The user has a working internet connection. \\
        {[D6]} & The external map service is available and provides accurate topographical data.\\
        & \\
        \hline
        & \\[-1em]
        
        % Requirements
        {[R1]} & The platform shall allow new users to register an account by providing personal information. \\
        {[R2]} & The platform shall allow registered users to log in by entering their registered email and password. \\
        {[R8]} & The platform shall allow users to start the navigation mode for the currently selected path. \\
        {[R12]} & When the navigation is completed, the platform shall present to registered users a summary screen
        containing statistics and the map of the path taken. \\
        {[R17]} & The platform shall display real-time metrics during the trip, including duration, current speed,
        distance traveled, and weather conditions. \\
        {[R18]} & The platform shall allow registered users to stop the recording and immediately transition to the summary screen to review the captured data. \\
        {[R25]} & The platform shall display a dashboard listing the user’s historical trips and their data, sorted chronologically.\\
        {[R29]} & The platform shall display a user profile screen summarizing the main aggregate statistics. \\
        & \\
        \hline
    \end{tabularx}
\end{table}

\newpage
\begin{table}[H]
    \centering
    % Aumenta lo spazio tra le righe per leggibilità
    \renewcommand{\arraystretch}{1.5}
    
    % Definizione tabella: larghezza totale testo, bordo verticale, colonna stretta per ID, colonna larga X per testo
    \begin{tabularx}{\textwidth}{|l X|}
        \hline
        % --- CORREZIONE QUI SOTTO ---
        % Uso 'p' con una larghezza calcolata (larghezza totale - margini interni) invece di 'l'
        % Questo obbliga il testo lungo ad andare a capo.
        \multicolumn{2}{|p{\dimexpr\textwidth-2\tabcolsep}|}{\textbf{G2}: When available, users would like to be informed about meteorological information related to the paths of interest.} \\
        \hline
        & \\[-1em] % Piccolo aggiustamento spaziatura
        
        % Domain Assumptions
        {[D2]} & The user has a working internet connection. \\
        {[D5]} & The external weather information service is reliable.\\
        & \\
        \hline
        & \\[-1em]
        
        % Requirements
        {[R1]} & The platform shall allow new users to register an account by providing personal information. \\
        {[R2]} & The platform shall allow registered users to log in by entering their registered email and password. \\
        {[R17]} & The platform shall display real-time metrics during the trip, including duration, current speed, distance traveled, and weather conditions. \\
        {[R25]} & The platform shall display a dashboard listing the user’s historical trips and their data, sorted chronologically. \\ % Corretto: rimosso '&' extra
        & \\
        \hline
    \end{tabularx}
\end{table}

%------------------------------------------------------------------------------------------------------------------------------------------------


\clearpage
{\color{Blue}{\section{Formal Analysis Using Alloy}}}
\label{sect:alloy}
\noindent
The objective of this section is to validate the critical aspects of the Best Bike Paths (BBP) system using Alloy, a lightweight formal method based on relational logic. In particular, the modeling effort concentrates on the robustness of the report submission mechanism, a core feature for maintaining data reliability within the platform.
\subsection{Signatures}
The signatures defined below represent the vocabulary used to describe the BBP ecosystem. Here, we outline the primary atoms of the model, including User, Trip, and Path, while also defining the dynamic elements such as Obstacle and PathStatus that are essential for tracking the system's evolution.

\lstinputlisting[language=alloy]{RASD/Alloy/Signatures.txt}


\newpage
\subsection{Facts}
The facts represent the constraints that the system must respect and are always true in the model. For the BBP application, we defined numerous facts to enforce data consistency, such as the necessity for a PathReport to be strictly linked to a Trip performed by the author. Furthermore, they dictate the logic for the dynamic evolution of the system, determining how PathStatus and Obstacles are updated based on the aggregation of user reports.
\lstinputlisting[language=alloy]{RASD/Alloy/Facts.txt}

\subsection{Assertions}
The assertions are the safety properties that we want to verify in the model. These conditions must hold true to ensure the platform avoids invalid or illogical states. For the BBP scenario, we verified several critical aspects, such as the strict requirement that a User must have actually completed a Trip on a specific Path before submitting a report (AuthorWasThere), as well as the constraint that limits users to a single report submission per time step to prevent spam.
\lstinputlisting[language=alloy]{RASD/Alloy/Assertions.txt}

\subsection{Model Visualization}
The following model was obtained with the following predicate:
\lstinputlisting[language=alloy]{RASD/Alloy/Predicate.txt}

\vspace{5mm}
\noindent
\textbf{Timeframe 0 (focus on discard trip):}

\begin{figure}[!h]
    \centering
    \includegraphics[width=1.1\linewidth]{RASD/Images/Alloy/Timeframe0.png}
    \label{fig:placeholder}
\end{figure}
\noindent
User3 takes a trip on Path1 but decides to discard it; therefore, no report is submitted. The status of all paths remains unchanged, and the set of obstacles remains the same.

\newpage
\textbf{Timeframe 1 (focus on tie and add/remove obstacles):}
\begin{figure}[!h]
    \centering
    \includegraphics[width=1.1\linewidth]{RASD/Images/Alloy/Timeframe1.png}
    \label{fig:placeholder}
\end{figure}

\noindent
User0 emits a report on Path0 with status = 'Manteinance Required'. \\
User1 emits a report on Path0 with status = 'Sufficient' and one obstacle (Obstacle3).\\
User2 emits a report on Path2 with status = 'Manteinance Required' and two obstacles (Obstacle1 and Obstacle4).\\
User3 emits a report on Path1 with status = 'Medium' and one obstacle (Obstacle 2), and it removes Obstacle0.\\
As a result, Path1 and Path2 take the status of the single report associated, while since there is a tie in the count of report statuses for Path0, it is assigned the lower of the two statuses (Maintenance Required).\\
In addition, Obstacle0 will be replaced with Obstacle3 in Path1, Obstacle1 and Obstacle4 have been added to Path2, and Obstacle2 has been added to Path1.

\newpage
\textbf{Timeframe 2 (focus on merging):}
\begin{figure}[!h]
    \centering
    \includegraphics[width=1.1\linewidth]{RASD/Images/Alloy/Timeframe2.png}
    \label{fig:placeholder}
\end{figure}
\\
User0 emits a report on Path2 with status = 'Medium'. \\
User1 emits a report on Path0 with status = 'Medium'.\\
User3 emits a report on Path2 with status = 'Sufficient'.\\
User2 makes a trip on Path2 but decides to discard it, therefore no report is submitted.\\
As a result, the status of Path2 is updated to Medium according to the majority vote, while the remaining paths are unaffected.

%------------------------------------------------------------------------------------------------------------------------------------------------
\clearpage
{\color{Blue}{\section{Use of Generative AI}}}
\label{sect:generativeai}
\subsection{AI-Based Support Tools}

While tools such as ChatGPT were utilized transversally to support drafting, code refinement, and logical requirement analysis, the specific phase of User Interface Prototyping relied on models specialized in image generation.

\noindent
Specifically, Google AI Studio was employed to interact with the Gemini model. Thanks to its multimodal capabilities, it was possible to translate the textual descriptions of Use Cases into high-fidelity visual mockups. This approach allowed for a rapid visualization of the spatial organization of elements (layout).

\subsection{Example of Interaction: Mockup Generation}

 A significant example of this interaction is represented by the definition of the Homepage. Initially, a prompt was provided to the model aimed at establishing the general structure and color palette, inspired by leading market competitors.

\begin{PromptBox}
\quad Generate a mockup for the homepage of a mobile GPS navigation application. The interface must be in "Dark Mode", featuring a dark night-time street map background similar to Google Maps. At the bottom of the screen, include a navigation bar with four sections: "Home", "Published", "Recent Trip", and "Profile". Position two round floating buttons stacked vertically immediately above this bar in the bottom right corner: one for "Manual Mode" and the other for "Automated Mode".

\vspace{0.3cm}

\end{PromptBox}

\begin{GeminiBad}
\quad Here is a proposal for the Homepage design in night mode. The generated image presents: \begin{itemize} \item A dark-toned cartographic background (midnight blue and charcoal) with highlighted roads. \item A bottom navigation bar featuring the four requested sections. \item The two action buttons ("Manual" and "Automated") positioned at the bottom right for easy thumb access. \end{itemize}

\textbf{} \\
    \begin{center}
    \includegraphics[height=6cm, keepaspectratio]{RASD/Images/WrongImage.png}
    \end{center}
    
\vspace{0.3cm}
\end{GeminiBad}

\noindent
Although the structure was correct, the first iteration lacked crucial elements for immediate usability, such as a clear indication of the user's position and the destination input field. A second refinement prompt was therefore necessary to integrate these details without altering the already approved visual style.

\vspace{0.3cm}
\begin{PromptBox}
\quad Based on the image just generated, modify only the top area and the center of the map. Add a floating search bar at the top of the screen specifically for entering a destination (with a placeholder like "Where to?"). Additionally, place a prominent cursor, such as a glowing blue arrow, in the center of the map to indicate the user's current position. Ensure that the rest of the existing interface, including the bottom navigation bar, remains unchanged.

\vspace{0.3cm}
\end{PromptBox}


\begin{GeminiGood}
\quad Updated the mockup by integrating the requested elements: \begin{itemize} \item \textbf{Search Bar:} Inserted at the top with rounded corners and a search icon, maintaining the transparency typical of modern UIs. \item \textbf{Position Cursor:} Added a blue indicator with a slight glow effect at the center of the map to ensure immediate contrast against the dark background. \end{itemize} The rest of the layout (navbar and mode buttons) remained unchanged.

\textbf{} \\
    \begin{center}
    \includegraphics[height=6cm, keepaspectratio]{RASD/Images/TrueImage.png}
    \end{center}
    
\vspace{0.3cm}

\end{GeminiGood}

\vspace{0.5cm}



\noindent
This two-phase process allowed for the separation of the structural layout definition (navbar, map, buttons) from the functional details (search bar, cursor). The model demonstrated the ability to maintain stylistic consistency between generations, integrating the requested modifications naturally within the pre-existing visual context.

\newpage
\subsection{Example of Interaction: Rules Identification}

The following interaction describes a consultation aimed at defining the Standard Compliance section. We engaged the AI to clarify the complex requirements related to the handling of continuous geolocation history and the integration of external map services. Through two prompts, we progressed from identifying the applicable external regulations (GDPR, Terms of Service) to defining the internal technical solutions required to ensure compliance without compromising the platform’s functionality.

\vspace{0.5cm}

\begin{PromptBox}
\quad I am developing a mobile app for cyclists that features real-time navigation and trip tracking, relying on third-party map services. Considering I will be handling continuous geolocation data, what are the primary privacy regulations and map licensing constraints I need to strictly follow?

\vspace{0.3cm}

\end{PromptBox}

\begin{ChatGPTBad}
\quad Based on your description, the platform operates at the intersection of sensitive user data and proprietary services. Here are the mandatory constraints:
\begin{enumerate}
    \item Privacy and Sensitive Data: Since the system handles Personally Identifiable Information (PII) and continuous geolocation history, this constitutes a significant privacy risk as it reveals user habits and private locations. You must strictly comply with the General Data Protection Regulation (GDPR) and potentially the California Consumer Privacy Act (CCPA).
    \item Map Provider Terms of Service (ToS):The dependency on external providers (e.g., OpenStreetMap or Google Maps) mandates rigorous adherence to their ToS. You are required to:
    \begin{itemize}
        \item Dynamically display proper copyright attributions.
        \item Strictly enforce technical limitations on offline data caching to avoid licensing violations.
    \end{itemize}
\end{enumerate}
    
\vspace{0.3cm}
\end{ChatGPTBad}



\vspace{0.5cm}

\noindent
The analysis of the regulatory landscape revealed a critical tension: the application requires continuous data collection for functionality, yet strictly limits data retention for compliance. Recognizing that simple adherence to rules is insufficient, we directed the next query toward defining the architectural drivers necessary to resolve this conflict at the structural level.

\vspace{0.3cm}
\newpage
\begin{PromptBox}
\quad Based on these regulatory constraints, how should I design the software architecture to ensure compliance? Specifically, I need to understand how to embed "Privacy-by-Design" principles into the core system and how to technically enforce the caching limitations, while keeping the architecture flexible enough to adapt to future changes.

\vspace{0.3cm}
\end{PromptBox}


\begin{ChatGPTGood}
\quad To ensure compliance without compromising functionality, the system must be architected with specific drivers:
\begin{itemize} \item \textbf{Privacy-by-Design Core:} You must embed privacy-by-design principles directly into the software core, rather than treating compliance as an external add-on. This minimizes the risk of exposing user habits. \item \textbf{Adaptability and Flexibility:} The system must be architected with the flexibility to adapt to this dual set of regulatory and licensing constraints. Regulations and ToS change frequently; a modular design will allow you to update compliance logic without rewriting the entire application.\item \textbf{Technical Enforcement:} The architecture must include a logic layer that automatically manages data retention, ensuring you strictly enforce technical limitations on offline data caching as required by the map providers.\end{itemize} The rest of the layout (navbar and mode buttons) remained unchanged.
    
\vspace{0.3cm}

\end{ChatGPTGood}


\vspace{0.5cm}

\noindent \textbf{Verification} \\
We verified the information generated by the AI by checking official online sources. The details regarding privacy regulations and map service policies were confirmed to be accurate and consistent with real-world requirements. Therefore, the content is reliable and has been validated for use in this project.

\subsection{Conclusion}
The integration of AI-based tools has significantly improved efficiency in the preparation of complex technical documents. Nevertheless, despite their advantages, it remains essential to adopt a critical approach and carefully review the produced content to ensure accuracy, coherence, and completeness. Achieving high-quality results often requires multiple iterations and precise refinements, guiding the process step by step. AI tools have proven to be valuable support instruments, but they should be used with careful supervision throughout all phases of document development.

%------------------------------------------------------------------------------------------------------------------------------------------------

\clearpage
{\color{Blue}{\section{Effort Spent}}}
\label{sect:effort}
\subsection*{Andrea Ferritti}

\begin{table}[H]
    \centering
\begin{tabular}{|l|c|}
        \hline
        \textbf{Section} & \textbf{Hours} \\ \hline
        Introduction & 10 \\ \hline
        Overall Description & 9\\ \hline
        Specific Requirements & 16\\ \hline
        Formal Analysis Using Alloy & 11 \\ \hline
        Use of Generative AI & 1 \\ \hline
    \end{tabular}
    \caption{Effort - Andrea Ferritti}
    \label{tab:effortferritti}
\end{table}

\subsection*{Jieke Lin}

\begin{table}[H]
    \centering
\begin{tabular}{|l|c|}
        \hline
        \textbf{Section} & \textbf{Hours} \\ \hline
        Introduction & 10 \\ \hline
        Overall Description & 8 \\ \hline
        Specific Requirements & 16\\ \hline
        Formal Analysis Using Alloy & 12 \\ \hline
        Use of Generative AI & 1 \\ \hline
    \end{tabular}
    \caption{Effort - Jieke Lin}
    \label{tab:effortlin}
\end{table}

\subsection*{Tommaso Tosca}

\begin{table}[H]
    \centering
\begin{tabular}{|l|c|}
        \hline
        \textbf{Section} & \textbf{Hours} \\ \hline
        Introduction & 9 \\ \hline
        Overall Description & 11\\ \hline
        Specific Requirements & 16\\ \hline
        Formal Analysis Using Alloy & 14 \\ \hline
        Use of Generative AI & 1 \\ \hline
    \end{tabular}
    \caption{Effort - Tommaso Tosca}
    \label{tab:efforttosca}
\end{table}


%------------------------------------------------------------------------------------------------------------------------------------------------




%\clearpage
%\addcontentsline{toc}{section}{References}
%\bibliographystyle{plain}
%\bibliography{main}
%------------------------------------------------------------------------------------------------------------------------------------------------




\end{document}
